%% toc: P3 39 39.8
\section{Proofs: existential pebble games and fixed-subgraph queries}
\label{sec:annex-p03-w03-p03-061-existential}

Conventions are those of \cref{def:p03-w03-p03-061-setting}; \(C\) abbreviates
\(C_{\mathbf A,\mathbf B}\) when the structures are clear. The truth of a
formula under an assignment depends only on the values of its free variables;
this is an induction on formulas whose infinitary clauses are immediate, and
it is used without mention.

\begin{annexitem}{annex:p03-w03-p03-061-substitution}{thm:p03-w03-p03-061-substitution}{Substitution of terms for variables}
\begin{proof}
For a term \(t\) over \(V_k\) and \(\bar a\in A^k\) write
\(\rho(\bar a)\defeq(\rho(1)^{\mathbf A}[\bar a],\ldots,
\rho(k)^{\mathbf A}[\bar a])\). Define \(\varphi_\rho\) by recursion on
\(\varphi\), simultaneously for all \(\rho\), and verify the equivalence and
the free-variable condition along the way.

\emph{Basic formulas.} \(\varphi_\rho\) replaces each variable \(v_j\) in
\(\varphi\) by the term \(\rho(j)\). Its free variables are those of the terms
\(\rho(j)\) with \(v_j\) in \(\varphi\). The value of the replaced term at
\(\bar a\) is the value of \(v_j\) at \(\rho(\bar a)\), and constants are
unchanged, so the equivalence holds. Atomic formulas go to atomic formulas,
and \(t\neq t'\) to a negated equality, so \(\varphi_\rho\) stays in
\(\mathcal L\).

\emph{Conjunction and disjunction.} Put
\((\bigwedge\Psi)_\rho\defeq\bigwedge\{\psi_\rho:\psi\in\Psi\}\) and likewise
for \(\bigvee\). The equivalence and the free-variable condition pass through
each conjunct or disjunct. A finite \(\Psi\) gives a finite set, so finitary
formulas stay finitary.

\emph{Existential quantifier.} Let \(\varphi=\exists v_j\,\psi\). The set
\(W\) of variables occurring in the terms \(\rho(i)\), \(i\neq j\), has at
most \(k-1\) members, so some \(v_l\in V_k-W\) exists; fix the least such
\(l\). Let \(\rho'\) agree with \(\rho\) except that \(\rho'(j)=v_l\), and
put \(\varphi_\rho\defeq\exists v_l\,\psi_{\rho'}\). For \(\bar a\in A^k\)
and \(a\in A\) let \(\bar a[l/a]\) replace the \(l\)-th entry by \(a\). For
\(i\neq j\) the term \(\rho(i)\) does not contain \(v_l\), so
\(\rho'(\bar a[l/a])\) agrees with \(\rho(\bar a)\) at every \(i\neq j\) and
has \(a\) at \(j\). Hence
\[
\begin{aligned}
 \mathbf A\models\exists v_l\,\psi_{\rho'}[\bar a]
 &\iff\exists a\in A\;\mathbf A\models\psi\bigl[\rho'(\bar a[l/a])\bigr]\\
 &\iff\exists a\in A\;\mathbf A\models\psi\bigl[\rho(\bar a)[j/a]\bigr]\\
 &\iff\mathbf A\models\exists v_j\,\psi\bigl[\rho(\bar a)\bigr],
\end{aligned}
\]
using the induction hypothesis for \(\psi\) and \(\rho'\) in the first step.
The free variables of \(\psi_{\rho'}\) lie among the variables of the terms
\(\rho'(i)\) with \(v_i\) free in \(\psi\). Those with \(i=j\) contribute only
\(v_l\), which is now bound; those with \(i\neq j\) are terms \(\rho(i)\)
with \(v_i\) free in \(\exists v_j\psi\). The variables of
\(\varphi_\rho\) lie in \(V_k\), and no negation or universal quantifier has
been introduced.
\end{proof}
\end{annexitem}

\begin{annexitem}{annex:p03-w03-p03-061-datalog}{thm:p03-w03-p03-061-datalog}{Bounded-variable Datalog is existential finite-variable}
\begin{proof}
\emph{The stages and the least fixed point.} The bodies contain the
intensional symbols only positively, so \(\Phi\) is monotone:
\(\bar P\subseteq\bar P'\) componentwise implies
\(\Phi(\bar P)\subseteq\Phi(\bar P')\). Since \(\Phi^0\) is empty,
\(\Phi^0\subseteq\Phi^1\), and monotonicity gives
\(\Phi^m\subseteq\Phi^{m+1}\) by induction. Let \(U\defeq\bigcup_m\Phi^m\).
Monotonicity gives \(\Phi^{m+1}=\Phi(\Phi^m)\subseteq\Phi(U)\) for all
\(m\), so \(U\subseteq\Phi(U)\). Conversely, a tuple of \(\Phi(U)\) comes from
a rule and an assignment under which its finitely many intensional body atoms
hold in \(U\); each holds in some stage, and as the stages increase all hold
in \(\Phi^m\) for one \(m\), so the tuple lies in \(\Phi^{m+1}\subseteq U\).
Thus \(U\) is a fixed point. If \(\Phi(\bar P)\subseteq\bar P\), induction
gives \(\Phi^m\subseteq\bar P\) for every \(m\), since
\(\Phi^{m+1}=\Phi(\Phi^m)\subseteq\Phi(\bar P)\subseteq\bar P\). Hence
\(U=\Phi^\infty\).

\emph{Item 1.} Induction on \(m\). Put \(\theta^0_{i,\mu}\defeq\bot\); this is
correct as \(\Phi^0_i=\emptyset\). Assume the formulas for stage \(m\) exist.
Fix \(i\), \(\mu\) and a rule \(r\) with head \(S_i(u_1,\ldots,u_n)\). Rename
the distinct body variables of \(r\) as \(v_1,\ldots,v_p\) (\(p\le k\)); the
head variables not in the body keep separate names. If some intensional body
atom has \(\theta^m=\bot\) (as computed below), \(r\) contributes nothing to
\(\Phi^{m+1}\) and is dropped. Otherwise define:
\begin{itemize}
\item \(\chi\): the conjunction of the body literals, where an extensional
      atom, equality or inequality is kept, and an intensional atom
      \(S_j(w_1,\ldots,w_q)\) is replaced by \(\theta^m_{j,\nu}\) with
      \(\nu(t)\defeq w_t\). By the induction hypothesis,
      \(\mathbf A\models\chi[\bar d]\) exactly when every body literal holds
      in \((\mathbf A,\Phi^m)\) under \(v_j\mapsto d_j\).
\item \(J\): the set of \(j\le p\) such that \(v_j\) occurs in the head; for
      \(j\in J\) fix a head position \(i_j\) with \(u_{i_j}=v_j\).
\item \(\chi'\defeq\exists v_{j_1}\cdots\exists v_{j_e}\,\chi\), where
      \(j_1,\ldots,j_e\) list \(\{1,\ldots,p\}-J\). Its free variables lie
      in \(\{v_j:j\in J\}\).
\item \(\rho\): \(\rho(j)\defeq\mu(i_j)\) for \(j\in J\), \(\rho(j)\defeq
      v_j\) otherwise; \(\chi'_\rho\) is given by
      \cref{thm:p03-w03-p03-061-substitution}.
\item \(\varepsilon\): the conjunction of the equalities
      \(\mu(i)=\mu(i')\) for head positions with \(u_i=u_{i'}\) a variable,
      and \(\mu(i)=c\) for head positions with \(u_i\) a constant \(c\).
\end{itemize}
Put \(\psi_{r}\defeq\varepsilon\wedge\chi'_\rho\) and let
\(\theta^{m+1}_{i,\mu}\) be the disjunction of the \(\psi_r\) over the
remaining rules with head symbol \(S_i\) (\(\bot\) if none remains). The
variables are in \(V_k\); only atomic formulas, inequalities where the
program has them, finite conjunctions and disjunctions and existential
quantifiers are used, so \(\theta^{m+1}_{i,\mu}\in\exists\mathrm{FO}^k\)
(respectively \(\exists\mathrm{FO}^k(\neq)\)) or is \(\bot\). By
\cref{thm:p03-w03-p03-061-substitution} the free variables of
\(\chi'_\rho\) lie among the terms \(\mu(i_j)\), and those of
\(\varepsilon\) among the \(\mu(i)\).

Write \(\bar e\defeq(\mu(1)^{\mathbf A}[\bar a],\ldots,\mu(n)^{\mathbf
A}[\bar a])\) and \(\bar c\defeq\rho(\bar a)\). It remains to show that
\(\mathbf A\models\psi_r[\bar a]\) exactly when some assignment \(\beta\) of
the variables of \(r\) makes the body true in \((\mathbf A,\Phi^m)\) and
\(u_i^{\mathbf A}[\beta]=e_i\) for every \(i\).

Suppose \(\mathbf A\models\psi_r[\bar a]\). By
\cref{thm:p03-w03-p03-061-substitution}, \(\mathbf A\models\chi'[\bar c]\),
so some \(\bar d\in A^k\) agrees with \(\bar c\) on \(J\) and satisfies
\(\chi\). Let \(\beta(v_j)\defeq d_j\) for \(j\le p\). For a head variable
\(x\) not in the body choose a position \(i\) with \(u_i=x\) and let
\(\beta(x)\defeq e_i\); by \(\varepsilon\) the choice of \(i\) does not
matter. The body holds under \(\beta\). For a head position \(i\): if
\(u_i=v_j\) with \(j\in J\), then
\(\beta(v_j)=d_j=c_j=\mu(i_j)^{\mathbf A}[\bar a]=e_{i_j}=e_i\), the last
step by \(\varepsilon\); if \(u_i\) is a head-only variable,
\(\beta(u_i)=e_i\) by \(\varepsilon\); if \(u_i=c\), then \(e_i=c^{\mathbf
A}\) by \(\varepsilon\).

Conversely, given such \(\beta\), the equalities of \(\varepsilon\) hold at
\(\bar a\), because positions carrying the same term receive the same value
\(u_i^{\mathbf A}[\beta]\). Let \(d_j\defeq\beta(v_j)\) for \(j\le p\) and
\(d_j\defeq a_j\) for \(j>p\). Then \(\mathbf A\models\chi[\bar d]\), and for
\(j\in J\), \(d_j=u_{i_j}^{\mathbf A}[\beta]=e_{i_j}=c_j\). So \(\bar d\)
witnesses \(\mathbf A\models\chi'[\bar c]\), that is,
\(\mathbf A\models\chi'_\rho[\bar a]\).

Taking the disjunction over the rules, \(\mathbf A\models
\theta^{m+1}_{i,\mu}[\bar a]\) exactly when \(\bar e\in\Phi(\Phi^m)_i=
\Phi^{m+1}_i\). Dropped rules contribute nothing, because their bodies
contain an atom false at stage \(m\).

\emph{Item 2.} \(\Phi^\infty_i=\bigcup_m\Phi^m_i\) by the first paragraph, so
the disjunction of the \(\theta^m_{i,\mu}\) is correct. It is a disjunction of
a countable set of formulas with variables in \(V_k\), hence in
\(\exists\mathrm L^k_{\infty\omega}\) (respectively
\(\exists\mathrm L^k_{\infty\omega}(\neq)\)).

\emph{Item 3.} If the goal symbol \(S_i\) has arity \(n\le k\), take
\(\mu(t)\defeq v_t\). Then \(\theta^\infty_{i,\mu}\) has free variables among
\(v_1,\ldots,v_n\) and holds of \((a_1,\ldots,a_n)\) exactly when this tuple
is in \(\Phi^\infty_i\), on every \(\sigma\)-structure. For \(n=0\) it is a
sentence.
\end{proof}
\end{annexitem}

\begin{annexitem}{annex:p03-w03-p03-061-strategy}{thm:p03-w03-p03-061-strategy}{Strategies and families}
\begin{proof}
The argument is the same for both games; ``partial homomorphism'' means
``one-to-one partial homomorphism'' for the \((\exists,\neq,k)\)-game.

\emph{A family gives a strategy.} Let \(\mathcal I\) be a winning family.
Every member contains \(C\), and \(\mathcal I\) is nonempty and closed under
subfunctions, so \(C\in\mathcal I\). Duplicator keeps the invariant
\(f_p\in\mathcal I\), which holds at the empty position. If Spoiler removes a
pebble, Duplicator removes hers; the new relation is contained in \(f_p\) and
contains \(C\), so it is in \(\mathcal I\). If Spoiler places a pebble \(j\)
that is not in use on \(a\), at most \(k-1\) pebbles are in use, so
\(|f_p-C|\le k-1\). The forth property gives \(g\in\mathcal I\) with
\(f_p\subseteq g\) and \(a\in\operatorname{dom}(g)\); Duplicator puts pebble
\(j\) on \(g(a)\). The new relation \(f_p\cup\{(a,g(a))\}\) lies between
\(C\) and \(g\), so it is in \(\mathcal I\). Every member of \(\mathcal I\) is
a partial homomorphism, so Spoiler never wins.

\emph{A strategy gives a family.} Fix a winning strategy of Duplicator and let
\(\mathcal I\) be the set of relations \(f_p\) at positions \(p\) reached in
finite plays that follow the strategy. The empty position gives
\(C\in\mathcal I\). Each \(f_p\) is a partial homomorphism, as otherwise
Spoiler would have won at \(p\). At most \(k\) pebbles are down, so
\(|f_p-C|\le k\). For closure under subfunctions let \(C\subseteq g\subseteq
f_p\). Let Spoiler continue the play from \(p\) by removing each pebble whose
pair is not in \(g\); Duplicator's answers are forced. The resulting
relation contains \(C\) and the pairs of the remaining pebbles, which are the
pairs of \(g-C\), so it equals \(g\). For the forth property let \(|f_p-C|<k\)
and \(a\in A\). If all \(k\) pebbles are in use, two of them carry the same
pair or one carries a pair of \(C\), since otherwise \(|f_p-C|=k\); Spoiler
removes such a pebble, which leaves the relation unchanged. He then places a
free pebble on \(a\), and the strategy's answer \(b\) yields a position with
relation \(f_p\cup\{(a,b)\}\in\mathcal I\).

\emph{Fewer pebbles.} If \(\mathcal I\) is a winning family for \(k\) pebbles
and \(j\le k\), let \(\mathcal I_j\defeq\{f\in\mathcal I:|f-C|\le j\}\). It
contains \(C\) and is closed under subfunctions. If \(f\in\mathcal I_j\) has
\(|f-C|<j\le k\) and \(a\in A\), the forth property of \(\mathcal I\) gives
\(g\supseteq f\) with \(a\in\operatorname{dom}(g)\); then
\(f\cup\{(a,g(a))\}\) lies between \(C\) and \(g\), so it is in \(\mathcal I\),
and it has at most \(j\) pairs outside \(C\).
\end{proof}
\end{annexitem}

\begin{annexitem}{annex:p03-w03-p03-061-game-logic}{thm:p03-w03-p03-061-game-logic}{Existential pebble games characterize existential \(k\)-variable preservation}
\begin{proof}
By \cref{thm:p03-w03-p03-061-strategy} it suffices to work with winning
families. The proof is written for \(\exists\mathrm L^k_{\infty\omega}\);
the version with inequalities differs only at the two places marked.

\emph{(2)\(\Rightarrow\)(1).} Let \(\mathcal I\) be a winning family. Claim:
for every \(\varphi\in\exists\mathrm L^k_{\infty\omega}\), every
\(f\in\mathcal I\) and all \(\bar a\in A^k\), \(\bar b\in B^k\) such that
\(a_i\in\operatorname{dom}(f)\) and \(b_i=f(a_i)\) whenever \(v_i\) is free in
\(\varphi\),
\(\mathbf A\models\varphi[\bar a]\) implies \(\mathbf B\models\varphi[\bar
b]\). Induction on \(\varphi\), for all \(f,\bar a,\bar b\) at once.

For a term \(t\) of \(\varphi\), the value \(t^{\mathbf A}[\bar a]\) lies in
\(\operatorname{dom}(f)\) and \(f(t^{\mathbf A}[\bar a])=t^{\mathbf B}[\bar
b]\): for a free variable by hypothesis, for a constant because \(C\subseteq
f\). An atom \(R(t_1,\ldots,t_r)\) true in \(\mathbf A\) is therefore true in
\(\mathbf B\) because \(f\) is a partial homomorphism. An equality true in
\(\mathbf A\) is true in \(\mathbf B\) because \(f\) is a function. (With
inequalities: \(t\neq t'\) true in \(\mathbf A\) is true in \(\mathbf B\)
because \(f\) is injective.) The infinitary connectives are immediate, since
the free variables of a conjunct or disjunct are free in the compound.

Let \(\varphi=\exists v_j\,\psi\) and \(\mathbf A\models\psi[\bar a[j/a]]\).
Let \(F\) be the set of \(i\) with \(v_i\) free in \(\varphi\); then
\(j\notin F\) and \(|F|\le k-1\). The function
\(g\defeq C\cup\{(a_i,f(a_i)):i\in F\}\) lies between \(C\) and \(f\), so
\(g\in\mathcal I\), and \(|g-C|\le k-1<k\). The forth property gives
\(h\in\mathcal I\) with \(g\subseteq h\) and \(a\in\operatorname{dom}(h)\).
For \(i\in F\), \(h(a_i)=g(a_i)=b_i\). The induction hypothesis for
\(\psi\), \(h\), \(\bar a[j/a]\) and \(\bar b[j/h(a)]\) gives
\(\mathbf B\models\psi[\bar b[j/h(a)]]\), so \(\mathbf B\models\varphi[\bar
b]\).

For a sentence \(\varphi\) apply the claim with \(f=C\in\mathcal I\): every
\(\exists\mathrm L^k_{\infty\omega}\)-sentence true in \(\mathbf A\) is true
in \(\mathbf B\).

\emph{(1)\(\Rightarrow\)(2), and the finite case.} Let \(\mathcal L\) be
\(\exists\mathrm L^k_{\infty\omega}\), or \(\exists\mathrm{FO}^k\) when
\(\mathbf B\) is finite, and assume every \(\mathcal L\)-sentence true in
\(\mathbf A\) is true in \(\mathbf B\). For \(r\le k\) and tuples
\(\bar a\in A^r\), \(\bar b\in B^r\) write
\((\bar a)\le_{\mathcal L}(\bar b)\) when every \(\mathcal L\)-formula with
free variables among \(v_1,\ldots,v_r\) true of \(\bar a\) in \(\mathbf A\) is
true of \(\bar b\) in \(\mathbf B\). Let \(\mathcal I\) be the set of
relations \(C\cup\{(a_1,b_1),\ldots,(a_r,b_r)\}\) for all such \(r\),
\(\bar a,\bar b\) with \((\bar a)\le_{\mathcal L}(\bar b)\).

(a) Each member \(f\) is a partial homomorphism. Every element of
\(\operatorname{dom}(f)\) is the value \(t^{\mathbf A}[\bar a]\) of a term
over \(v_1,\ldots,v_r\), and \(f\) is the set of pairs
\((t^{\mathbf A}[\bar a],t^{\mathbf B}[\bar b])\). If
\(t^{\mathbf A}[\bar a]=t'^{\mathbf A}[\bar a]\), the atomic formula
\(t=t'\) holds of \(\bar a\), hence of \(\bar b\); so \(f\) is a function.
Relations are preserved in the same way, through atomic formulas. (With
inequalities: \(t\neq t'\) transfers in the same way, so \(f\) is
injective.) Also \(|f-C|\le r\le k\).

(b) \(C\in\mathcal I\): this is the case \(r=0\), which is the hypothesis.

(c) Re-indexing. If \((\bar a)\le_{\mathcal L}(\bar b)\) and
\(1\le i_1,\ldots,i_q\le r\), then \((a_{i_1},\ldots,a_{i_q})
\le_{\mathcal L}(b_{i_1},\ldots,b_{i_q})\). Indeed, let \(\varphi\) have free
variables among \(v_1,\ldots,v_q\) and let \(\rho(t)\defeq v_{i_t}\) for
\(t\le q\), \(\rho(t)\defeq v_t\) otherwise. By
\cref{thm:p03-w03-p03-061-substitution}, \(\varphi_\rho\in\mathcal L\) has
free variables among \(v_1,\ldots,v_r\) and holds of \(\bar a\) exactly
when \(\varphi\) holds of \((a_{i_1},\ldots,a_{i_q})\), and likewise in
\(\mathbf B\). Consequently \(\mathcal I\) is closed under subfunctions: if
\(C\subseteq g\subseteq f\) with \(f\) given by \(\bar a,\bar b\), list the
pairs of \(g-C\) as \((a_{i_t},b_{i_t})\). It also follows that every
\(f\in\mathcal I\) has a representation by tuples of length
\(|f-C|\).

(d) Forth. Let \(f\in\mathcal I\), \(|f-C|<k\) and \(a\in A\). By (c),
\(f\) is given by \(\bar a,\bar b\) of length \(r=|f-C|<k\). Suppose that
for every \(b\in B\) the relation \((\bar a,a)\le_{\mathcal L}(\bar b,b)\)
fails, and choose \(\varphi_b\in\mathcal L\) with free variables among
\(v_1,\ldots,v_{r+1}\), true of \((\bar a,a)\) and false of \((\bar b,b)\).
Put \(\psi\defeq\exists v_{r+1}\bigwedge_{b\in B}\varphi_b\). Its variables
lie in \(V_k\), and it is in \(\mathcal L\): when \(\mathcal L\) is
first-order, \(B\) is finite and the conjunction is finitary. \(\psi\) holds
of \(\bar a\) in \(\mathbf A\), with witness \(a\), so it holds of \(\bar b\)
in \(\mathbf B\) with some witness \(b\), and then \(\varphi_b\) holds of
\((\bar b,b)\), a contradiction. Hence some \(b\) has
\((\bar a,a)\le_{\mathcal L}(\bar b,b)\), and
\(f\cup\{(a,b)\}\in\mathcal I\) extends \(f\) with \(a\) in its domain.

So \(\mathcal I\) is a winning family. With \(\mathcal L=\exists\mathrm
L^k_{\infty\omega}\) this proves (1)\(\Rightarrow\)(2). If \(\mathbf B\) is
finite, \(\mathbf A\preceq^{\exists,k}_{\omega\omega}\mathbf B\) gives (2) by
the case \(\mathcal L=\exists\mathrm{FO}^k\), and (1) implies
\(\mathbf A\preceq^{\exists,k}_{\omega\omega}\mathbf B\) because
\(\exists\mathrm{FO}^k\subseteq\exists\mathrm L^k_{\infty\omega}\).
\end{proof}
\end{annexitem}

\begin{annexitem}{annex:p03-w03-p03-061-homomorphism}{thm:p03-w03-p03-061-homomorphism}{Preservation under homomorphisms}
\begin{proof}
A homomorphism of \(\sigma\)-structures sends \(c^{\mathbf A}\) to
\(c^{\mathbf B}\), so \(C\subseteq h\). Let \(\mathcal I\defeq\{f:C\subseteq
f\subseteq h,\ |f-C|\le k\}\). Its members are restrictions of \(h\), hence
partial homomorphisms (one-to-one if \(h\) is); it contains \(C\) and is
closed under subfunctions; and if \(|f-C|<k\), then
\(f\cup\{(a,h(a))\}\in\mathcal I\). By \cref{thm:p03-w03-p03-061-strategy}
Duplicator wins, and by \cref{thm:p03-w03-p03-061-game-logic} every
\(\exists\mathrm L^k_{\infty\omega}\)-sentence (respectively
\(\exists\mathrm L^k_{\infty\omega}(\neq)\)-sentence) true in \(\mathbf A\) is
true in \(\mathbf B\), for every \(k\).
\end{proof}
\end{annexitem}

\begin{annexitem}{annex:p03-w03-p03-061-cliques}{thm:p03-w03-p03-061-cliques}{The existential game is asymmetric}
\begin{proof}
The vocabulary is \(\{E\}\), so \(C=\emptyset\). Spoiler places his
\(k+1\) pebbles on the \(k+1\) distinct vertices \(a_1,\ldots,a_{k+1}\) of
\(\mathbf K_{k+1}\). Duplicator's answers lie in a set of \(k\) vertices, so
\(b_i=b_j\) for some \(i\neq j\). Either the relation \(f_p\) is not a
function, or it is and \(E(a_i,a_j)\) holds while \(E(b_i,b_i)\) fails, since
\(\mathbf K_k\) has no loops. Either way Spoiler wins.

For the other direction, any injection \(h\) of the vertices of \(\mathbf
K_k\) into those of \(\mathbf K_{k+1}\) is a one-to-one homomorphism, since
distinct vertices go to distinct, hence adjacent, vertices.
\Cref{thm:p03-w03-p03-061-homomorphism} gives a winning family for the
\((\exists,\neq,j)\)-game for every \(j\). A one-to-one partial homomorphism
is a partial homomorphism, so the same family wins the \((\exists,j)\)-game.
\end{proof}
\end{annexitem}

\begin{annexitem}{annex:p03-w03-p03-061-criterion}{thm:p03-w03-p03-061-criterion}{Game criterion for existential finite-variable definability}
\begin{proof}
It suffices to show that, for each fixed \(k\), item~2 for \(k\) holds
exactly when \(Q\) is \(\exists\mathrm L^k_{\infty\omega}\)-definable on
\(\mathcal C\); the equivalence of items 1 and 2 follows because
\(\exists\mathrm L^\omega_{\infty\omega}\) is the union of the
\(\exists\mathrm L^k_{\infty\omega}\).

Suppose \(\varphi\in\exists\mathrm L^k_{\infty\omega}\) defines \(Q\) on
\(\mathcal C\). If \(Q(\mathbf A)=1\) and Duplicator wins the
\((\exists,k)\)-game on \(\mathbf A,\mathbf B\in\mathcal C\), then
\(\mathbf A\models\varphi\), so \(\mathbf B\models\varphi\) by
\cref{thm:p03-w03-p03-061-game-logic}, and \(Q(\mathbf B)=1\).

Conversely assume item~2 for \(k\). Every finite structure is isomorphic to
one whose universe is \(\{1,\ldots,n\}\) for some \(n\), and these form a set;
let \(\mathcal R\) be the set of such structures that are isomorphic to a
member of \(\mathcal C\) with \(Q=1\). For \(\mathbf A\in\mathcal R\) let
\(\tau_{\mathbf A}\) be the conjunction of the set of
\(\exists\mathrm{FO}^k\)-sentences true in \(\mathbf A\), and let
\(\varphi\defeq\bigvee_{\mathbf A\in\mathcal R}\tau_{\mathbf A}\), a sentence
of \(\exists\mathrm L^k_{\infty\omega}\). If \(\mathbf B\in\mathcal C\) and
\(Q(\mathbf B)=1\), then \(\mathbf B\) is isomorphic to some \(\mathbf
A\in\mathcal R\), and \(\mathbf B\models\tau_{\mathbf A}\) because
isomorphisms preserve truth. If \(\mathbf B\in\mathcal C\) and
\(\mathbf B\models\tau_{\mathbf A}\) for some \(\mathbf A\in\mathcal R\), then
\(\mathbf A\preceq^{\exists,k}_{\omega\omega}\mathbf B\). Since \(\mathbf B\)
is finite, \cref{thm:p03-w03-p03-061-game-logic} shows that Duplicator wins the
\((\exists,k)\)-game on \(\mathbf A'\) and \(\mathbf B\), where
\(\mathbf A'\in\mathcal C\) is isomorphic to \(\mathbf A\) with
\(Q(\mathbf A')=1\): isomorphic structures satisfy the same sentences, so
\(\mathbf A'\preceq^{\exists,k}_{\omega\omega}\mathbf B\) as well. Item~2
gives \(Q(\mathbf B)=1\). Thus \(\varphi\) defines \(Q\) on \(\mathcal C\).

The version with inequalities is word for word the same, using
\(\exists\mathrm{FO}^k(\neq)\), \(\exists\mathrm L^k_{\infty\omega}(\neq)\)
and the \((\exists,\neq,k)\)-games.
\end{proof}
\end{annexitem}

\begin{annexitem}{annex:p03-w03-p03-061-horn}{thm:p03-w03-p03-061-horn}{Satisfiability as an existential pebble game}
\begin{proof}
An empty clause makes the conjunction unsatisfiable, just as
\(x\wedge\neg x\) does. Adding an unused variable changes no clause, and
any satisfying assignment extends to it. Thus the normalization in
\cref{def:p03-w03-p03-061-horn} preserves satisfiability and ensures a
nonempty encoding. In this proof write \(\varphi\) for the normalized
conjunction and allow assignments on its enlarged variable set.

\emph{(1)\(\Rightarrow\)(2).} A satisfying assignment \(s\) of \(\varphi\) is
a homomorphism \(\mathbf A_\varphi\to\mathbf B\). For each tuple of a relation
of \(\mathbf A_\varphi\), the corresponding clause is true under \(s\), and
the relation of \(\mathbf B\) with the same name consists exactly of the
value tuples satisfying that clause form. \Cref{thm:p03-w03-p03-061-homomorphism}
gives (2). \emph{(2)\(\Rightarrow\)(3)} is the case \(k=3\).

\emph{(3)\(\Rightarrow\)(1).} Let \(\pi\) be the program displayed after the
theorem, and on a \(\sigma\)-structure let \(T^\infty\) and the truth value
of the goal \(P\) be those of its least fixed point.

First, in \(\mathbf A_\varphi\), \(T^\infty\) is the least set \(M\) of
variables closed under the rules ``\(x\in M\) if \(x\) is a clause; \(z\in M\) if \(\neg
x\vee z\) is a clause and \(x\in M\); \(z\in M\) if \(\neg x\vee\neg y\vee z\)
is a clause and \(x,y\in M\)'', and these are the three rules for \(T\).
\(P\) holds in \(\mathbf A_\varphi\) exactly when some clause with no positive
literal has all its variables in \(M\).

Second, \(P\) holds in \(\mathbf A_\varphi\) exactly when \(\varphi\) is
unsatisfiable. If a negative clause has all its variables in \(M\): every
satisfying assignment \(s\) of \(\varphi\) makes all of \(M\) true (induction
along the closure rules, since each rule's clause forces its head true when
its premises are true), hence falsifies that clause, a contradiction. If no
negative clause lies within \(M\): the assignment true exactly on \(M\)
satisfies every clause. A positive unit clause \(x\) has \(x\in M\). A clause
\(\neg x\vee z\) or \(\neg x\vee\neg y\vee z\) is true unless its premises are
in \(M\), and then \(z\in M\). A negative clause has a variable outside \(M\).

Third, \(P\) fails in \(\mathbf B\). The rules for \(T\) give \(1\in T\) from
\(P_1^{\mathbf B}\). If \(x=1\) and \((x,z)\in P_2^{\mathbf B}\) then
\(z=1\), and if \(x=y=1\) and \((x,y,z)\in P_3^{\mathbf B}\) then \(z=1\), so
\(T^\infty=\{1\}\). The goal rules would need \(1\in N_1^{\mathbf B}\),
\((1,1)\in N_2^{\mathbf B}\) or \((1,1,1)\in N_3^{\mathbf B}\), and none of
these holds.

The program \(\pi\) is in \(3\)-Datalog: every body has at most three
variables and every head at most one. By
\cref{thm:p03-w03-p03-061-datalog} its goal is defined by an
\(\exists\mathrm L^3_{\infty\omega}\)-sentence. If Duplicator wins the
\((\exists,3)\)-game on \(\mathbf A_\varphi\) and \(\mathbf B\),
\cref{thm:p03-w03-p03-061-game-logic} gives
\(\mathbf A_\varphi\preceq^{\exists,3}_{\infty\omega}\mathbf B\). Since \(P\)
fails in \(\mathbf B\), it fails in \(\mathbf A_\varphi\), and \(\varphi\) is
satisfiable by the second step.
\end{proof}
\end{annexitem}

\begin{annexitem}{annex:p03-w03-p03-061-menger}{thm:p03-w03-p03-061-menger}{Menger's theorem for directed graphs, fan form}
\begin{proof}
Loops lie on no path and are deleted first. Call a path from \(X\) to \(Y\)
an \emph{\(X\)--\(Y\) path} when its first node is its only node in \(X\) and
its last node is its only node in \(Y\); a one-node path at a node of
\(X\cap Y\) is one. Every path from \(X\) to \(Y\) contains an
\(X\)--\(Y\) path (from its last node in \(X\) before its first node in
\(Y\) to that node), so a set meets every path from \(X\) to \(Y\) exactly
when it meets every \(X\)--\(Y\) path; call such a set
\emph{separating}. Disjoint paths from \(X\) to \(Y\) yield disjoint
\(X\)--\(Y\) paths.

\emph{Item 1.} Disjoint \(X\)--\(Y\) paths need distinct nodes of any
separating set, so their number is at most the least size \(\kappa\) of a
separating set. For the converse, induction on the number of edges. With no
edges, the \(X\)--\(Y\) paths are the one-node paths at \(X\cap Y\);
\(X\cap Y\) is separating and every separating set contains it, so
\(\kappa=|X\cap Y|\) and the one-node paths realize it.

Let \(e=(x,y)\) be an edge, \(x\neq y\). If \(\mathbf D-e\) has no separating
set of size less than \(\kappa\), induction gives \(\kappa\) disjoint paths in
\(\mathbf D-e\), which are paths in \(\mathbf D\). Otherwise let \(W\) be
separating in \(\mathbf D-e\) with \(|W|<\kappa\). Every path from \(X\) to
\(Y\) in \(\mathbf D\) that avoids \(W\) uses \(e\), so \(W_x\defeq
W\cup\{x\}\) and \(W_y\defeq W\cup\{y\}\) are separating in \(\mathbf D\).
Hence \(|W|=\kappa-1\) and \(x,y\notin W\).

Every set \(Z\) separating \(X\) from \(W_x\) in \(\mathbf D-e\) separates
\(X\) from \(Y\) in \(\mathbf D\). Let \(P\) be a path from \(X\) to \(Y\) in
\(\mathbf D\). It meets \(W_x\); its initial segment up to its first node in
\(W_x\) does not use \(e\), because if \(P\) uses \(e\) it reaches \(x\)
before \(y\). That segment is a path from \(X\) to \(W_x\) in \(\mathbf D-e\),
so it meets \(Z\). Thus \(|Z|\ge\kappa\), and by induction \(\mathbf D-e\) has
\(\kappa=|W_x|\) disjoint \(X\)--\(W_x\) paths, one ending at each node of
\(W_x\). Symmetrically, using final segments from the last node in \(W_y\),
which follow \(y\) when \(P\) uses \(e\), \(\mathbf D-e\) has \(\kappa\)
disjoint \(W_y\)--\(Y\) paths, one starting at each node of \(W_y\).

Let \(P\) be one of the first paths and \(R\) one of the second, and suppose
they share a node \(z\notin W\). If \(z\) is neither the last node of \(P\)
nor the first node of \(R\), then \(P\) up to \(z\) followed by \(R\) from
\(z\) is a walk in \(\mathbf D-e\) from \(X\) to \(Y\) avoiding \(W\) (the
first part meets \(W_x\) only in its last node, the second meets \(W_y\) only
in its first node), and it contains a path from \(X\) to \(Y\) in \(\mathbf
D-e\) avoiding \(W\), which is impossible. If \(z\) is the last node of
\(P\), then \(z=x\) and \(z\) is not the first node of \(R\) (that is in
\(W_y\) and \(x\notin W_y\)); the same walk again avoids \(W\). If \(z\) is
the first node of \(R\), then \(z=y\), and symmetrically. Hence the two
systems meet only in \(W\). Each node \(w\in W\) is the last node of exactly
one \(P_w\) and the first node of exactly one \(R_w\); \(P_w\) meets \(W\)
only in \(w\) and \(R_w\) meets \(W\) only in \(w\), so \(P_wR_w\) is a path.
Together with \(P_x\,e\,R_y\) (where \(P_x\) ends at \(x\) and \(R_y\) starts
at \(y\); they are disjoint, as \(P_x\) avoids \(W\)) these are \(\kappa\)
paths from \(X\) to \(Y\) in \(\mathbf D\). Two of them can meet only in
\(W\), and at a node \(w\) of \(W\) only \(P_wR_w\) passes, so they are
disjoint.

\emph{Item 2.} Apply item 1 to \(\mathbf D-s\) with \(X\) the set of
out-neighbours of \(s\) other than \(s\) and \(Y\defeq T\). If there are \(m\)
disjoint \(X\)--\(T\) paths \(R_1,\ldots,R_m\) in \(\mathbf D-s\), the paths
\(sR_i\) start at \(s\), meet \(T\) only in their last nodes and share only
\(s\). Otherwise item 1 gives a separating set \(Z\subseteq V-\{s\}\) with
\(|Z|\le m-1\). A walk from \(s\) to a node of \(T\) contains a path from
\(s\) to that node; as \(s\notin T\), the path has a second node, which is in
\(X\), and its part from there is a path from \(X\) to \(T\) in \(\mathbf
D-s\). It meets \(Z\).
\end{proof}
\end{annexitem}

\begin{annexitem}{annex:p03-w03-p03-061-star}{thm:p03-w03-p03-061-star}{Star patterns are existentially definable with inequalities}
\begin{proof}
\emph{Out-stars.} Let the centre be \(r\) and the leaves \(w_1,\ldots,w_m\).
Write \(s,t_i\) for \(d_r^{\mathbf G},d_{w_i}^{\mathbf G}\) and
\(s',t_i'\) for the corresponding nodes of \(\mathbf G'\), and
\(T'\defeq\{t'_1,\ldots,t'_m\}\). Since the \(t_i\) are distinct and differ
from \(s\), two paths \(p_i\) (from \(s\) to \(t_i\)) and \(p_j\) are
node-disjoint exactly when they share only \(s\), and then neither passes
through another leaf. So \(\mathrm{FiSH}(\mathbf H)(\mathbf G)=1\) gives
simple paths \(p_1,\ldots,p_m\) of \(\mathbf G\), \(p_i\) from \(s\) to
\(t_i\), pairwise sharing only \(s\). If \(\mathbf G'\) had \(m\) paths from
\(s'\) to \(T'\), each meeting \(T'\) only in its last node and any two
sharing only \(s'\), their last nodes would be distinct and so would be the
\(m\) nodes \(t'_i\); they would witness
\(\mathrm{FiSH}(\mathbf H)(\mathbf G')=1\). Hence by
\cref{thm:p03-w03-p03-061-menger} there is \(Z\subseteq V(\mathbf
G')-\{s'\}\) with \(|Z|\le m-1\) meeting every walk from \(s'\) to a node of
\(T'\).

Suppose a winning family \(\mathcal I\) for the \((\exists,\neq,m+1)\)-game on
\(\mathbf G\) and \(\mathbf G'\) existed. We build \(f_0,\ldots,f_m\in
\mathcal I\) and nodes \(x_1,\ldots,x_m\) of \(\mathbf G\) such that
\(f_i-C\) consists of pairs \((x_j,\cdot)\) with \(j\le i\), each \(x_j\) is a
node of \(p_j\) other than \(s\), \(x_j\in\operatorname{dom}(f_i)\), and
\(f_i(x_j)\in Z\) for \(j\le i\). Start with \(f_0\defeq C\in\mathcal I\).

Given \(f_{i-1}\) with \(i\le m\), let \(p_i=(u_0,u_1,\ldots,u_L)\) with
\(u_0=s\), \(u_L=t_i\), \(L\ge1\). Put \(h_0\defeq f_{i-1}\). For
\(l=1,\ldots,L\), given \(h_{l-1}\in\mathcal I\) with
\(h_{l-1}\subseteq f_{i-1}\cup\{(u_{l-1},h_{l-1}(u_{l-1}))\}\) and
\(u_{l-1}\in\operatorname{dom}(h_{l-1})\): \(|h_{l-1}-C|\le i\le m<m+1\), so
the forth property gives \(h'\in\mathcal I\) with \(h_{l-1}\subseteq h'\)
and \(u_l\in\operatorname{dom}(h')\). Let \(h_l\defeq f_{i-1}\cup
\{(u_l,h'(u_l))\}\), which lies between \(C\) and \(h'\) and so is in
\(\mathcal I\). Put \(y_l\defeq h'(u_l)\) and \(y_0\defeq s'\). Since
\(E^{\mathbf G}(u_{l-1},u_l)\) and \(h'\) is a partial homomorphism containing
\(h_{l-1}\), \(E^{\mathbf G'}(y_{l-1},y_l)\). As \(t_i\) is a constant,
\(y_L=t'_i\). So \(y_0,\ldots,y_L\) is a walk in \(\mathbf G'\) from \(s'\)
to \(t'_i\), and it meets \(Z\); let \(l^*\) be least with \(y_{l^*}\in Z\).
Then \(l^*\ge1\) because \(s'\notin Z\). Put \(x_i\defeq u_{l^*}\) and
\(f_i\defeq h_{l^*}\). The pair \((x_i,y_{l^*})\) is the only new pair, and
\(f_i(x_j)=f_{i-1}(x_j)\in Z\) for \(j<i\).

The nodes \(x_1,\ldots,x_m\) are distinct, since \(x_j\) lies on \(p_j\) and
differs from \(s\), and the paths share only \(s\). Their images under the
one-to-one function \(f_m\) are \(m\) distinct elements of \(Z\), which has
at most \(m-1\) elements. This contradiction shows that no winning family
exists, so by \cref{thm:p03-w03-p03-061-strategy} Spoiler wins the
\((\exists,\neq,m+1)\)-game, and by its last clause the
\((\exists,\neq,k)\)-game for every \(k\ge m+1\).

\emph{In-stars.} Reversing all edges turns an in-star \(\mathbf H\) into an
out-star \(\mathbf H^{\mathrm{rev}}\), and
\(\mathrm{FiSH}(\mathbf H)(\mathbf G)=\mathrm{FiSH}(\mathbf H^{\mathrm
{rev}})(\mathbf G^{\mathrm{rev}})\), since reversed simple paths are simple
paths with the same nodes. A function is a (one-to-one) partial homomorphism
from \(\mathbf G\) to \(\mathbf G'\) exactly when it is one from
\(\mathbf G^{\mathrm{rev}}\) to \(\mathbf G'^{\mathrm{rev}}\), so winning
families for the two pairs coincide, and the out-star case applies.

\emph{Definability.} By the first part, if
\(\mathrm{FiSH}(\mathbf H)(\mathbf G)=1\) and Duplicator wins the
\((\exists,\neq,m+1)\)-game on \(\mathbf G\) and \(\mathbf G'\) in
\(\mathcal F_{\mathbf H}\), then \(\mathrm{FiSH}(\mathbf H)(\mathbf G')=1\).
The query is invariant under isomorphism, so the fixed-\(k\) form of
\cref{thm:p03-w03-p03-061-criterion} with inequalities, for \(k=m+1\), shows
that it is \(\exists\mathrm L^{m+1}_{\infty\omega}(\neq)\)-definable on
\(\mathcal F_{\mathbf H}\).
\end{proof}
\end{annexitem}

\begin{annexitem}{annex:p03-w03-p03-061-fish-preserved}{thm:p03-w03-p03-061-fish-preserved}{Homeomorphism queries are preserved by one-to-one homomorphisms}
\begin{proof}
Let paths \(p_e\) witness \(\mathrm{FiSH}(\mathbf H)(\mathbf G)=1\). A
homomorphism sends edges to edges, so the image \(g(p_e)\) of each path is a
walk from \(g(d_u^{\mathbf G})=d_u^{\mathbf G'}\) to
\(d_w^{\mathbf G'}\) for \(e=(u,w)\). Since \(g\) is injective it repeats no
node, so it is a simple path. If \(g(p_e)\) and \(g(p_{e'})\) share a node
\(g(z)\), injectivity puts \(z\) on both \(p_e\) and \(p_{e'}\), so \(z\) is
an endpoint of both, and \(g(z)\) is an endpoint of both images. The images
witness \(\mathrm{FiSH}(\mathbf H)(\mathbf G')=1\).
\end{proof}
\end{annexitem}

\begin{annexitem}{annex:p03-w03-p03-061-star-datalog}{thm:p03-w03-p03-061-star-datalog}{Star patterns are Datalog\((\neq)\)-definable}
\begin{proof}
Write \(\Phi\) for the operator of \(\pi_m\) on \(\mathbf G\) and
\(R_j^n\), \(R_j^\infty\) for the components of its stages and of its least
fixed point. The first paragraph of the proof of
\cref{thm:p03-w03-p03-061-datalog} applies to every Datalog\((\neq)\)
program: the stages increase, their union is \(\Phi^\infty\), and
\(\Phi(\Phi^\infty)=\Phi^\infty\), so \(\Phi^\infty\) is closed under every
rule.

\emph{Walk description.} Fix \(j\), \(l=m-j\). For a tuple
\(\bar c=(s,y_1,\ldots,y_{j-1},x,t_1,\ldots,t_l)\) of nodes, say that a walk
\(s=w_0,w_1,\ldots,w_n=x\) with \(n\ge1\) is \emph{admissible} for
\(\bar c\) when \(s\) and every \(w_i\) with \(i\ge1\) differ from
\(t_1,\ldots,t_l\) and, if \(j\ge2\),
\[
 (s,y_1,\ldots,y_{j-1},w_i,t_1,\ldots,t_l)\in R_{j-1}^\infty
 \qquad(i\ge1).
\] Claim: \(\bar c\in R_j^\infty\) exactly when some walk is
admissible for \(\bar c\).

If an admissible walk of length \(n\) exists, induction on \(n\) shows
\(\bar c\in R_j^\infty\). For \(n=1\) the first rule for \(R_j\) applies,
since \(E(s,x)\) holds, the inequalities hold and the atom \(\Lambda_j\) is
in \(\Phi^\infty\); closure of \(\Phi^\infty\) under the rule gives the
tuple. For \(n>1\), the walk \(w_0,\ldots,w_{n-1}\) is admissible for the
tuple with \(w_{n-1}\) in place of \(x\), which is in \(R_j^\infty\) by
induction; the second rule with \(w\defeq w_{n-1}\) gives \(\bar c\).

Conversely, induction on the stage \(n\) shows that every tuple of \(R_j^n\)
has an admissible walk. \(R_j^0\) is empty. A tuple enters \(R_j^{n+1}\)
through a rule whose body holds in \(\Phi^n\). By the first rule, \(E(s,x)\)
holds, \(s\) and \(x\) differ from the \(t\)'s, and for \(j\ge2\) the atom
\(R_{j-1}(s,\bar y,x,\bar t)\) holds in \(\Phi^n\subseteq\Phi^\infty\); the
walk \(s,x\) is admissible. By the second rule, the tuple with \(w\) in place
of \(x\) is in \(R_j^n\) and has an admissible walk ending at \(w\) by
induction; appending the edge \((w,x)\), with \(x\) different from the
\(t\)'s and \(R_{j-1}(s,\bar y,x,\bar t)\) in \(\Phi^\infty\), gives an
admissible walk ending at \(x\).

\emph{The equivalence.} Induction on \(j\). Fix pairwise distinct
\(s,s_1,\ldots,s_j\) and nodes \(t_1,\ldots,t_l\), and let \(\bar c\) be the
tuple with \(y_i\defeq s_i\) (\(i<j\)) and \(x\defeq s_j\). For \(j\ge2\),
the tuples \((s,s_1,\ldots,s_{j-1},w,t_1,\ldots,t_l)\) have pairwise distinct
\(s,s_1,\ldots,s_{j-1}\), so by the induction hypothesis, for
\(l+1=m-(j-1)\), they lie in \(R_{j-1}^\infty\) exactly when
\(Q_{j-1,l+1}(s,s_1,\ldots,s_{j-1},w,t_1,\ldots,t_l)\) holds.

Suppose \(Q_{j,l}(s,s_1,\ldots,s_j,\bar t)\), witnessed by
\(p_1,\ldots,p_j\). Take the walk \(p_j=(w_0,\ldots,w_n)\); \(n\ge1\) since
\(s\neq s_j\). Its nodes, \(s\) included, differ from the \(t\)'s. For
\(j\ge2\) and \(i\ge1\), \(w_i\neq s\) because \(p_j\) is simple, so \(w_i\)
lies on none of \(p_1,\ldots,p_{j-1}\), which share with \(p_j\) only \(s\).
Hence \(p_1,\ldots,p_{j-1}\) witness
\(Q_{j-1,l+1}(s,s_1,\ldots,s_{j-1},w_i,\bar t)\). The walk is admissible,
and \(\bar c\in R_j^\infty\).

Suppose \(\bar c\in R_j^\infty\), with admissible walk
\(W=(w_0,\ldots,w_n)\), and suppose \(Q_{j,l}(s,s_1,\ldots,s_j,\bar t)\)
fails. For \(j=1\), \(W\) is a walk from \(s\) to \(s_1\) avoiding the
\(t\)'s; it contains a simple path from \(s\) to \(s_1\) with the same
property, which witnesses \(Q_{1,l}\). So \(j\ge2\). The nodes
\(s_1,\ldots,s_{j-1}\) differ from the \(t\)'s, since they are endpoints of
the paths witnessing \(Q_{j-1,l+1}\) at \(w_n\); \(s\) and \(s_j=w_n\) differ
from them by admissibility. Let \(\mathbf D\) be the subgraph of
\(\mathbf G\) induced on the nodes other than \(t_1,\ldots,t_l\), and
\(T\defeq\{s_1,\ldots,s_j\}\), a set of \(j\) nodes of \(\mathbf D\) not
containing \(s\). If \(\mathbf D\) had \(j\) paths from \(s\) to \(T\), each
meeting \(T\) only in its last node and any two sharing only \(s\), their
last nodes would be distinct, hence all of \(T\), and they would witness
\(Q_{j,l}\). By \cref{thm:p03-w03-p03-061-menger}, item~2, there is
\(Z\subseteq V(\mathbf D)-\{s\}\) with \(|Z|\le j-1\) meeting every walk of
\(\mathbf D\) from \(s\) to a node of \(T\). The walk \(W\) lies in
\(\mathbf D\) and ends at \(s_j\in T\), so \(w_i\in Z\) for some \(i\ge1\).
By admissibility and the induction hypothesis there are paths
\(p_1,\ldots,p_{j-1}\), \(p_r\) from \(s\) to \(s_r\), any two sharing only
\(s\), avoiding \(w_i\) and the \(t\)'s. Each lies in \(\mathbf D\), so each
meets \(Z\), in a node other than \(s\) and other than \(w_i\). Distinct
paths meet \(Z\) in distinct nodes, since they share only \(s\). So
\(Z-\{w_i\}\), of size at most \(j-2\), has \(j-1\) elements, a
contradiction.

\emph{The goal and the variable count.} Let \(\mathbf H\) be an out-star
with centre \(r\) and leaves \(w_1,\ldots,w_m\), \(\mathbf G\in\mathcal
F_{\mathbf H}\), \(s\defeq d_r^{\mathbf G}\), \(s_i\defeq d_{w_i}^{\mathbf
G}\). These nodes are pairwise distinct, and simple paths \(p_i\) from \(s\)
to \(s_i\) are node-disjoint exactly when any two share only \(s\), because
\(s\) is the only endpoint that \(p_i\) and \(p_{i'}\) have in common. So
\(\mathrm{FiSH}(\mathbf H)(\mathbf G)=1\) exactly when
\(Q_{m,0}(s,s_1,\ldots,s_m)\), exactly when \((s,s_1,\ldots,s_m)\in
R_m^\infty\), exactly when the goal holds. For an in-star, the program with
reversed atoms computes on \(\mathbf G\) what \(\pi_m\) computes on the
reversed graph, and the reversal argument of the proof of
\cref{thm:p03-w03-p03-061-star} applies. Every body of \(\pi_m\) has at most
the variables \(s,y_1,\ldots,y_{j-1},w,x,t_1,\ldots,t_{m-j}\), that is
\(m+2\) of them, and every head at most \(m+1\); the goal rule has none.
\end{proof}
\end{annexitem}

\begin{annexitem}{annex:p03-w03-p03-061-basic}{thm:p03-w03-p03-061-basic}{Every non-star pattern contains a basic pattern}
\begin{proof}
Let \(\mathbf H\) be a pattern graph that is not a star graph. Patterns have
no loops, so the two nodes of an edge are distinct.

If two edges have no common node, their four nodes are distinct and form a
copy of \(\mathbf H_1\). If some node \(v\) is the head of an edge
\((u,v)\) and the tail of an edge \((v,w)\), then \(u\neq v\neq w\); for
\(u=w\) the two edges form a copy of \(\mathbf H_3\), and otherwise a copy
of \(\mathbf H_2\).

Suppose neither case occurs: any two edges share a node, and no node is both
a head and a tail. \(\mathbf H\) has an edge \((a,b)\). If every edge had
tail \(a\), every node other than \(a\) would, not being isolated, be the
head of an edge from \(a\), and \(\mathbf H\) would be an out-star with
centre \(a\). So some edge \((c,d)\) has \(c\neq a\); similarly, as
\(\mathbf H\) is not an in-star with centre \(b\), some edge \((c',d')\) has
\(d'\neq b\). The edge \((c,d)\) shares a node with \((a,b)\). Now \(c=b\)
would make \(b\) a head and a tail, and \(d=a\) would do so for \(a\);
hence \(d=b\) and \(c\notin\{a,b\}\). In the same way \((c',d')\) shares a
node with \((a,b)\), \(d'=a\) and \(c'=b\) are excluded, and so \(c'=a\) and
\(d'\notin\{a,b\}\). The edges \((c,b)\) and \((a,d')\) share a node; since
\(c\neq a\), \(b\neq a\) and \(b\neq d'\), this forces \(c=d'\). Then \(c\)
is the head of \((a,c)\) and the tail of \((c,b)\), a contradiction.
\end{proof}
\end{annexitem}

\begin{annexitem}{annex:p03-w03-p03-061-transfer}{thm:p03-w03-p03-061-transfer}{Non-definability passes to larger patterns}
\begin{proof}
\emph{Reduction to a subgraph.} Let \(\iota\) be an isomorphism from
\(\mathbf F_1\) onto a subgraph \(\mathbf F_1'\) of \(\mathbf F_2\).
Renaming each constant symbol \(d_v\) as \(d_{\iota(v)}\) is a bijection
between \(\sigma_{\mathbf F_1}\)-structures and
\(\sigma_{\mathbf F_1'}\)-structures with the same graph. A homeomorphism of
\(\mathbf F_1\) extending \(v\mapsto d_v\) is, through \(\iota\), a
homeomorphism of \(\mathbf F_1'\) extending \(\iota(v)\mapsto
d_{\iota(v)}\) with the same paths. The renaming carries formulas of
\(\exists\mathrm L^k_{\infty\omega}(\neq)\) to formulas of the same logic
with the same truth values. So \(\mathrm{FiSH}(\mathbf F_1')\) is not
\(\exists\mathrm L^k_{\infty\omega}(\neq)\)-definable either, and we may
assume that \(\mathbf F_1=(V_1,E_1)\) is a subgraph of
\(\mathbf F_2=(V_2,E_2)\).

\emph{Distinguished nodes on paths.} Let \(\mathbf F\) be a pattern graph,
\(\mathbf G\in\mathcal F_{\mathbf F}\), and let paths \(p_e\)
(\(e\in E_{\mathbf F}\)) form a homeomorphism extending
\(v\mapsto d_v^{\mathbf G}\). Then for \(e=(u,w)\) the path \(p_e\) contains
no node \(d_x^{\mathbf G}\) with \(x\notin\{u,w\}\). Indeed, \(x\) is not
isolated, so it lies on an edge \(e'\), and \(e'\neq e\). The path
\(p_{e'}\) has \(d_x^{\mathbf G}\) as an endpoint, and \(p_e,p_{e'}\) are
node-disjoint, so \(d_x^{\mathbf G}\) would be an endpoint of \(p_e\), that
is \(d_u^{\mathbf G}\) or \(d_w^{\mathbf G}\); the constants denote distinct
nodes, so \(x\in\{u,w\}\), a contradiction.

\emph{The construction.} By the fixed-\(k\) form of
\cref{thm:p03-w03-p03-061-criterion} with inequalities (the query is
invariant under isomorphism), there are \(\mathbf A,\mathbf B\in\mathcal
F_{\mathbf F_1}\) with \(\mathrm{FiSH}(\mathbf F_1)(\mathbf A)=1\),
\(\mathrm{FiSH}(\mathbf F_1)(\mathbf B)=0\), such that Duplicator wins the
\((\exists,\neq,k)\)-pebble game on \(\mathbf A\) and \(\mathbf B\); let
\(\mathcal I\) be a winning family (\cref{thm:p03-w03-p03-061-strategy}).
Let \(\mathbf A'\) have universe \(A\) together with new elements \(a'_v\)
for \(v\in V_2-V_1\); let \(d_v^{\mathbf A'}\defeq d_v^{\mathbf A}\) for
\(v\in V_1\) and \(d_v^{\mathbf A'}\defeq a'_v\) otherwise; and let
\(E^{\mathbf A'}\defeq E^{\mathbf A}\cup\{(d_u^{\mathbf A'},d_w^{\mathbf
A'}):(u,w)\in E_2-E_1\}\). Build \(\mathbf B'\) from \(\mathbf B\) in the same
way with new elements \(b'_v\). The constants of \(\mathbf A'\) and of
\(\mathbf B'\) are pairwise distinct, so both lie in
\(\mathcal F_{\mathbf F_2}\).

\emph{\(\mathrm{FiSH}(\mathbf F_2)(\mathbf A')=1\).} Let \(p_e\)
(\(e\in E_1\)) witness \(\mathrm{FiSH}(\mathbf F_1)(\mathbf A)=1\); they are
paths of \(\mathbf A'\). For \(e=(u,w)\in E_2-E_1\) let \(q_e\) be the
one-edge path from \(d_u^{\mathbf A'}\) to \(d_w^{\mathbf A'}\), simple
because \(u\neq w\). A node shared by \(q_e\) and another path of the family
is an endpoint of \(q_e\), say \(d_x^{\mathbf A'}\). If the other path is
some \(q_{e'}\), the node is an endpoint of it too. If it is some
\(p_{e'}\), then \(x\in V_1\) (the new elements are not in \(A\)) and, by
the previous paragraph applied to \(\mathbf F_1\) and \(\mathbf A\), \(x\) is
an endpoint of \(e'\), so the node is an endpoint of \(p_{e'}\). Two paths
\(p_e,p_{e'}\) are node-disjoint by hypothesis. Hence the paths \(p_e\) and
\(q_e\) are pairwise node-disjoint and witness
\(\mathrm{FiSH}(\mathbf F_2)(\mathbf A')=1\).

\emph{\(\mathrm{FiSH}(\mathbf F_2)(\mathbf B')=0\).} Suppose paths \(p_e\)
(\(e\in E_2\)) witness the contrary. Let \(e=(u,w)\in E_1\). By the
paragraph on distinguished nodes, applied to \(\mathbf F_2\) and
\(\mathbf B'\), the only distinguished nodes on \(p_e\) are its endpoints
\(d_u^{\mathbf B'}\) and \(d_w^{\mathbf B'}\). The new elements are
distinguished, so \(p_e\) runs in \(B\). An edge of
\(E^{\mathbf B'}-E^{\mathbf B}\) joins two distinguished nodes, so on
\(p_e\) it would be \((d_u^{\mathbf B'},d_w^{\mathbf B'})\), since no edge of
a simple path leaves its last node or enters its first; but such an edge is
\((d_x^{\mathbf B'},d_y^{\mathbf B'})\) for some \((x,y)\in E_2-E_1\), and the
constants are distinct, so \((x,y)=(u,w)\in E_1\), which is impossible. So
\(p_e\) is a path of \(\mathbf B\). The constants of \(\mathbf F_1\) have the
same values in \(\mathbf B\) and \(\mathbf B'\), so the paths \(p_e\)
(\(e\in E_1\)) witness \(\mathrm{FiSH}(\mathbf F_1)(\mathbf B)=1\), a
contradiction.

\emph{Duplicator wins on \(\mathbf A'\) and \(\mathbf B'\).} Let
\(N\defeq\{(a'_v,b'_v):v\in V_2-V_1\}\), so that
\(C_{\mathbf A',\mathbf B'}=C_{\mathbf A,\mathbf B}\cup N\), and
\(\mathcal I'\defeq\{f\cup N:f\in\mathcal I\}\). Each \(f\cup N\) is a
one-to-one function, as \(f\) and \(N\) are and their domains and ranges are
disjoint. It sends \(d_v^{\mathbf A'}\) to \(d_v^{\mathbf B'}\) for every
\(v\in V_2\). An edge of \(\mathbf A\) between elements of its domain goes to
an edge of \(\mathbf B\); an added edge \((d_u^{\mathbf A'},d_w^{\mathbf
A'})\) goes to the added edge \((d_u^{\mathbf B'},d_w^{\mathbf B'})\). So
\(f\cup N\) is a one-to-one partial homomorphism, and
\(|(f\cup N)-C_{\mathbf A',\mathbf B'}|=|f-C_{\mathbf A,\mathbf B}|\le k\).
If \(C_{\mathbf A',\mathbf B'}\subseteq g\subseteq f\cup N\), then
\(g=g_0\cup N\) with \(C_{\mathbf A,\mathbf B}\subseteq g_0\subseteq f\), so
\(g_0\in\mathcal I\) and \(g\in\mathcal I'\). For the forth property, a new
element is already in every domain, and for \(a\in A\) the forth property of
\(\mathcal I\) gives \(g\supseteq f\) with \(a\in\operatorname{dom}(g)\), and
then \(g\cup N\supseteq f\cup N\). So \(\mathcal I'\) is a winning family.

By \cref{thm:p03-w03-p03-061-strategy} and the fixed-\(k\) form of
\cref{thm:p03-w03-p03-061-criterion}, the pair \(\mathbf A',\mathbf B'\)
shows that \(\mathrm{FiSH}(\mathbf F_2)\) is not
\(\exists\mathrm L^k_{\infty\omega}(\neq)\)-definable on
\(\mathcal F_{\mathbf F_2}\).

\emph{The consequence.} Let \(\mathbf H\) be a non-star pattern and suppose
\(\mathrm{FiSH}(\mathbf H)\) is defined by a sentence of
\(\exists\mathrm L^k_{\infty\omega}(\neq)\). By
\cref{thm:p03-w03-p03-061-basic}, some \(\mathbf H_i\) is isomorphic to a
subgraph of \(\mathbf H\), and by the first part
\(\mathrm{FiSH}(\mathbf H_i)\) is
\(\exists\mathrm L^k_{\infty\omega}(\neq)\)-definable, contrary to the
hypothesis.
\end{proof}
\end{annexitem}

\begin{annexitem}{annex:p03-w03-p03-061-walks}{thm:p03-w03-p03-061-walks}{Walks of every fixed length with three variables}
\begin{proof}
Induction on \(m\). The formula \(\psi_1=E(v_1,v_2)\) is atomic, its free
variables are \(v_1,v_2\), and it holds at \((a_1,a_2,a_3)\) exactly when
\(E^{\mathbf A}(a_1,a_2)\), that is, when \(e_0=a_1,e_1=a_2\) is a walk of
length \(1\).

Assume the claim for \(\psi_m\). The variables of \(\psi_{m+1}\) are those
of \(\psi_m\) together with \(v_1,v_3\), so they lie in \(\{v_1,v_2,v_3\}\);
only atomic formulas, conjunction and existential quantifiers are used, so
\(\psi_{m+1}\in\exists\mathrm{FO}^3\). In
\((\exists v_1)(v_1=v_3\wedge\psi_m)\) the variable \(v_1\) is bound, so its
free variables lie among \(v_2,v_3\); in \(E(v_1,v_3)\wedge(\exists
v_1)(v_1=v_3\wedge\psi_m)\) they lie among \(v_1,v_2,v_3\), and
\((\exists v_3)\) binds \(v_3\). So the free variables of \(\psi_{m+1}\) lie
among \(v_1,v_2\).

For \(a_1,a_2,a_3\in A\), \(\mathbf A\models\psi_{m+1}[a_1,a_2,a_3]\) holds
exactly when there is \(c\in A\) with \(E^{\mathbf A}(a_1,c)\) and some
\(d\in A\) with \(d=c\) and \(\mathbf A\models\psi_m[d,a_2,c]\), that is,
with \(\mathbf A\models\psi_m[c,a_2,c]\). By the induction hypothesis the last
condition says that a walk \(c=e_1,\ldots,e_{m+1}=a_2\) of length \(m\)
exists. Prefixing \(e_0\defeq a_1\) gives a walk of length \(m+1\) from
\(a_1\) to \(a_2\); conversely, the second node \(c\defeq e_1\) of such a
walk satisfies \(E^{\mathbf A}(a_1,c)\) and starts a walk of length \(m\) to
\(a_2\).
\end{proof}
\end{annexitem}

\begin{annexitem}{annex:p03-w03-p03-061-system}{thm:p03-w03-p03-061-system}{Stages and least fixed points of positive existential systems}
\begin{proof}
The proof is written for \(\exists\mathrm{FO}^k\); with inequalities nothing
changes, since a negated equality is replaced by itself below and does not
mention the \(S_j\).

\emph{Removing \(\bot\).} Call a formula \emph{extended} if it is built from
\(\bot\) and \(\exists\mathrm{FO}^k\)-formulas by finite conjunction, finite
disjunction and existential quantification over variables of \(V_k\).
Every extended formula \(\chi\) is equivalent to \(\bot\) or to an
\(\exists\mathrm{FO}^k\)-formula \(\chi^\circ\) whose free variables are among
those of \(\chi\). Induction on \(\chi\): a conjunction with a conjunct
equivalent to \(\bot\) is equivalent to \(\bot\), and otherwise the
conjunction of the \(\chi_i^\circ\) is used; a disjunction drops its
disjuncts equivalent to \(\bot\) and is \(\bot\) if none remains;
\(\exists v_j\,\chi\) is equivalent to \(\bot\) if \(\chi\) is, and to
\(\exists v_j\,\chi^\circ\) otherwise.

\emph{Item 1.} Induction on \(m\); \(\theta^0_i\defeq\bot\) defines
\(\Phi^0_i=\emptyset\). Given \(\theta^m_1,\ldots,\theta^m_l\), let
\(\varphi_i^*\) arise from \(\varphi_i\) by replacing each atom
\(S_j(t_1,\ldots,t_{n_j})\) by \(\bot\) if \(\theta^m_j=\bot\), and otherwise
by \((\theta^m_j)_\rho\) of \cref{thm:p03-w03-p03-061-substitution}, with
\(\rho(r)\defeq t_r\) for \(r\le n_j\) and \(\rho(r)\defeq v_r\) for
\(n_j<r\le k\). The free variables of \((\theta^m_j)_\rho\) lie among the
terms \(\rho(r)\) with \(v_r\) free in \(\theta^m_j\), hence among the
\(t_r\) with \(r\le n_j\), which are terms of the replaced atom. So
\(\varphi_i^*\) is an extended formula with variables in \(V_k\) and free
variables among those of \(\varphi_i\); let \(\theta^{m+1}_i\) be \(\bot\) or
\((\varphi_i^*)^\circ\) accordingly.

Claim: for every subformula \(\psi\) of \(\varphi_i\), its replacement
\(\psi^*\) and all \(\bar a\in A^k\),
\(\mathbf A\models\psi^*[\bar a]\iff(\mathbf A,\Phi^m)\models\psi[\bar a]\).
For atoms of \(\sigma\), equalities and negated equalities nothing is
replaced. For an atom \(S_j(t_1,\ldots,t_{n_j})\) with \(\theta^m_j\neq\bot\),
\cref{thm:p03-w03-p03-061-substitution} gives
\(\mathbf A\models(\theta^m_j)_\rho[\bar a]\iff\mathbf
A\models\theta^m_j[\rho(\bar a)]\), and as the free variables of
\(\theta^m_j\) are among \(v_1,\ldots,v_{n_j}\), the induction hypothesis
turns this into \((t_1^{\mathbf A}[\bar a],\ldots,t_{n_j}^{\mathbf A}[\bar
a])\in\Phi^m_j\), which is the truth of the atom in \((\mathbf A,\Phi^m)\);
if \(\theta^m_j=\bot\), both sides are false since \(\Phi^m_j=\emptyset\).
Replacement commutes with conjunction, disjunction and \(\exists v_j\), so
the claim passes to compound formulas. For \(\psi=\varphi_i\) it says that
\(\theta^{m+1}_i\) defines \(\Phi(\Phi^m)_i=\Phi^{m+1}_i\).

\emph{Item 2.} \emph{Monotonicity.} If \(\bar P\subseteq\bar P'\)
componentwise, then \((\mathbf A,\bar P)\models\psi[\bar a]\) implies
\((\mathbf A,\bar P')\models\psi[\bar a]\) for every negation-free formula
\(\psi\), by induction on \(\psi\): atoms \(S_j\) pass from \(P_j\) to
\(P'_j\), the other basic formulas do not mention the \(S_j\), and
conjunction, disjunction and \(\exists\) preserve the implication. Hence
\(\Phi(\bar P)\subseteq\Phi(\bar P')\).

\emph{Least fixed point.} By monotonicity and \(\Phi^0\subseteq\Phi^1\), the
stages increase. Let \(U\defeq\bigcup_m\Phi^m\). Monotonicity gives
\(\Phi^{m+1}\subseteq\Phi(U)\), so \(U\subseteq\Phi(U)\). For the converse,
by induction on a finitary negation-free formula \(\psi\): if \((\mathbf
A,U)\models\psi[\bar a]\), then \((\mathbf A,\Phi^m)\models\psi[\bar a]\) for
some \(m\). An atom \(S_j\) true in \(U\) is true in some stage; the other
basic formulas do not mention the \(S_j\); for a finite conjunction take the
largest of the finitely many stages and use that the stages increase; for a
disjunction or an existential formula use the true disjunct or the witness.
Applied to \(\varphi_i\), this puts every tuple of \(\Phi(U)_i\) into some
\(\Phi^{m+1}_i\subseteq U_i\), so \(\Phi(U)\subseteq U\). If
\(\Phi(\bar P)\subseteq\bar P\), induction gives \(\Phi^m\subseteq\bar P\)
for all \(m\), so \(U\subseteq\bar P\). Hence \(U\) is the least fixed point.

\emph{Definability.} A tuple lies in \(\Phi^\infty_i\) exactly when it lies
in some \(\Phi^m_i\), exactly when some \(\theta^m_i\) holds of it. The
disjunction of the countable set \(\{\theta^m_i:m\ge0\}\) has its variables
in \(V_k\) and uses no negation, so it belongs to
\(\exists\mathrm L^k_{\infty\omega}\).
\end{proof}
\end{annexitem}

\begin{annexitem}{annex:p03-w03-p03-061-proper}{thm:p03-w03-p03-061-proper}{The inclusion of bounded-variable Datalog is proper}
\begin{proof}
\emph{Item 1.} Let \(\pi\) be a Datalog\((\neq)\) program and \(\mathbf A\) a
finite \(\sigma\)-structure. Each rule has finitely many variables, so the
assignments of its variables into \(A\) can be listed, and from the list of
the tuples of a stage \(\Phi^m\) the stage \(\Phi^{m+1}\) is computed by
checking every rule under every assignment. The stages increase (first
paragraph of the proof of \cref{thm:p03-w03-p03-061-datalog}) within the
finite set of tuples of the intensional relations, so
\(\Phi^{m+1}=\Phi^m\) for some \(m\) at most the number of those tuples.
Then \(\Phi^{m'}=\Phi^m\) for all \(m'\ge m\), and
\(\Phi^\infty=\bigcup_{m'}\Phi^{m'}=\Phi^m\). The algorithm computes stages
until two consecutive ones agree and reads off the goal.

\emph{Item 2.} With \(\psi_n\) of \cref{thm:p03-w03-p03-061-walks}, put
\[
 \chi_S\defeq\bigvee_{p\in S}(\exists v_1)(\exists v_2)\bigl(v_1=v_2\wedge
 \psi_p\bigr),
\]
a sentence of \(\exists\mathrm L^3_{\infty\omega}\). By
\cref{thm:p03-w03-p03-061-walks} the disjunct for \(p\) holds in \(\mathbf A\)
exactly when some walk of length \(p\) leads from a node to itself, so
\(\chi_S\) defines \(Q_S\).

For a prime \(p\) let \(\mathbf C_p\) be the \(\sigma\)-structure with universe
\(\{0,\ldots,p-1\}\), \(E^{\mathbf C_p}=\{(r,r+1\bmod p):r<p\}\), all other
relations empty and all constants equal to \(0\). A walk in
\(\mathbf C_p\) advances by \(1\) modulo \(p\) at each step, so a closed walk
has a length divisible by \(p\), and the walk once around the cycle has
length \(p\). The only prime divisible by \(p\) is \(p\), so
\(Q_S(\mathbf C_p)=1\) exactly when \(p\in S\).

There are uncountably many sets of primes, as there are infinitely many
primes, and only countably many algorithms. Choose \(S\) such that no
algorithm decides membership in \(S\) on the primes. If an algorithm
decided \(Q_S\), then on input a prime \(p\) one could write down
\(\mathbf C_p\) and run it, deciding \(p\in S\); so \(Q_S\) is not
computable.

\emph{Consequences.} Fix \(k\ge1\). By \cref{thm:p03-w03-p03-061-datalog},
every Boolean \(k\)-Datalog query is
\(\exists\mathrm L^k_{\infty\omega}\)-definable, hence
\(\exists\mathrm L^\omega_{\infty\omega}\)-definable, and the same holds with
inequalities. The query \(Q_S\) is
\(\exists\mathrm L^3_{\infty\omega}\)-definable, hence definable in
\(\exists\mathrm L^\omega_{\infty\omega}\) and in
\(\exists\mathrm L^\omega_{\infty\omega}(\neq)\), and for \(k\ge3\) in
\(\exists\mathrm L^k_{\infty\omega}\) and
\(\exists\mathrm L^k_{\infty\omega}(\neq)\), since \(V_3\subseteq V_k\). By
item~1 no Datalog\((\neq)\) program defines it. So each inclusion is proper.
\end{proof}
\end{annexitem}

\begin{annexitem}{annex:p03-w03-p03-061-spoiler-lfp}{thm:p03-w03-p03-061-spoiler-lfp}{Determining the winner of the existential game}
\begin{proof}
The proof is written for the \((\exists,k)\)-pebble game; for the
\((\exists,\neq,k)\)-pebble game read ``one-to-one partial homomorphism''
throughout, and the additions are marked. Write \(C\) for
\(C_{\mathbf A,\mathbf B}\), \(f_{\bar a\bar b}\) for \(f_{\bar a,\bar b}\)
and \(\mathcal W\) for \(\mathcal W_{\mathbf A,\mathbf B}\). In the
definition of \(\Gamma\), \(T\) occurs only in the condition
``\((\bar a[j/a],\bar b[j/b])\in T\) for every \(b\)'', which is preserved
when \(T\) grows; so \(\Gamma\) is monotone, and by
\cref{thm:p03-w03-p03-059-knaster-tarski} its least fixed point exists and is
the intersection of all \(T\) with \(\Gamma(T)\subseteq T\).

\emph{Item 1, first direction.} Let \(\mathcal I\) be a winning family and
\(J\defeq\{(\bar a,\bar b):f_{\bar a\bar b}\in\mathcal I\}\), \(X\) its
complement in \(A^k\times B^k\). We show \(\Gamma(X)\subseteq X\), that is,
no \((\bar a,\bar b)\in J\) lies in \(\Gamma(X)\). Its relation
\(f_{\bar a\bar b}\) is in \(\mathcal I\), hence a partial homomorphism. Let
\(j\le k\) and \(a\in A\). The relation \(g\defeq C\cup\{(a_i,b_i):i\neq j\}\)
lies between \(C\) and \(f_{\bar a\bar b}\), so \(g\in\mathcal I\), and
\(|g-C|\le k-1\). The forth property gives \(h\in\mathcal I\) with
\(g\subseteq h\) and \(a\in\operatorname{dom}(h)\); put \(b\defeq h(a)\).
Then \(f_{\bar a[j/a]\bar b[j/b]}=g\cup\{(a,b)\}\) lies between \(C\) and
\(h\), so it is in \(\mathcal I\) and \((\bar a[j/a],\bar b[j/b])\notin X\).
So neither condition of \(\Gamma(X)\) holds at \((\bar a,\bar b)\). By the
intersection formula \(\mathcal W\subseteq X\), that is, \(\mathcal W\cap
J=\emptyset\).

Now let \(\bar a\in A^k\). Put \(g_0\defeq C\in\mathcal I\). Given
\(g_{i-1}=C\cup\{(a_r,b_r):r<i\}\in\mathcal I\) with \(i\le k\), we have
\(|g_{i-1}-C|\le i-1<k\); the forth property gives \(h\in\mathcal I\)
containing \(g_{i-1}\) with \(a_i\) in its domain, and with
\(b_i\defeq h(a_i)\) the relation \(g_i\defeq g_{i-1}\cup\{(a_i,b_i)\}\) lies
between \(C\) and \(h\), so \(g_i\in\mathcal I\). Then
\(f_{\bar a\bar b}=g_k\in\mathcal I\), so \((\bar a,\bar b)\in J\) and
\((\bar a,\bar b)\notin\mathcal W\).

\emph{Item 1, second direction.} Assume that every \(\bar a\) has some
\(\bar b\) with \((\bar a,\bar b)\notin\mathcal W\). Let \(\mathcal I\) be the
set of relations \(g\) with \(C\subseteq g\subseteq f_{\bar a\bar b}\) for
some \((\bar a,\bar b)\notin\mathcal W\). As \(\mathcal W=\Gamma(\mathcal W)\),
such \((\bar a,\bar b)\) is not in \(\Gamma(\mathcal W)\), so
\(f_{\bar a\bar b}\) is a partial homomorphism, and so is every
\(g\subseteq f_{\bar a\bar b}\) containing \(C\). Also
\(|g-C|\le|f_{\bar a\bar b}-C|\le k\). \(A\) is nonempty, so \(\mathcal I\)
is nonempty, and it is closed under subfunctions by definition. Forth: let
\(g\in\mathcal I\) with \(C\subseteq g\subseteq f_{\bar a\bar b}\),
\((\bar a,\bar b)\notin\mathcal W\), \(|g-C|<k\), and \(a\in A\). Each pair
of \(g-C\) is \((a_i,b_i)\) for some \(i\); choose one such \(i\) for each
pair, giving a set \(I_0\) of fewer than \(k\) indices, and fix \(j\notin
I_0\). Since \((\bar a,\bar b)\notin\Gamma(\mathcal W)\), some \(b\in B\) has
\((\bar a[j/a],\bar b[j/b])\notin\mathcal W\). Then
\(g\cup\{(a,b)\}\subseteq C\cup\{(a_i,b_i):i\in I_0\}\cup\{(a,b)\}
\subseteq f_{\bar a[j/a]\bar b[j/b]}\), so \(g\cup\{(a,b)\}\in\mathcal I\).
Hence \(\mathcal I\) is a winning family, and Duplicator wins by
\cref{thm:p03-w03-p03-061-strategy}.

\emph{Item 2.} Let \(\bar x=x_1,\ldots,x_k\) and \(\bar y=y_1,\ldots,y_k\) be
distinct variables. For a term \(t\) over \(\bar x\) and the constant symbols
let \(t^1\) replace \(x_i\) by itself and \(c\) by \(c^1\), and \(t^2\)
replace \(x_i\) by \(y_i\) and \(c\) by \(c^2\). Let \(\vartheta(\bar x,\bar
y)\) be the disjunction of the formulas \(t^1=t'^1\wedge\neg\,t^2=t'^2\) for
all terms \(t,t'\), and
\(R^1(t_1^1,\ldots,t_r^1)\wedge\neg R^2(t_1^2,\ldots,t_r^2)\) for all
relation symbols \(R\) and terms \(t_1,\ldots,t_r\) (with inequalities add
\(\neg\,t^1=t'^1\wedge t^2=t'^2\)). The elements of
\(f_{\bar a\bar b}\) are the pairs \((t^{\mathbf A}[\bar a],t^{\mathbf
B}[\bar b])\), so \(\mathbf A\oplus\mathbf B\models\vartheta[\bar a,\bar b]\)
exactly when \(f_{\bar a\bar b}\) is not a function, or (with inequalities)
not injective, or is a function that fails to preserve some relation, that
is, when it is not a (one-to-one) partial homomorphism. Let
\[
\begin{aligned}
 \psi(T,\bar x,\bar y)\defeq{}&\bigwedge_{i\le k}\bigl(D_1(x_i)\wedge
 D_2(y_i)\bigr)\wedge\Bigl(\vartheta\\
 &\vee\bigvee_{j\le k}(\exists x_j)\bigl(
 D_1(x_j)\wedge(\forall y_j)(\neg D_2(y_j)\vee T(\bar x,\bar y))\bigr)\Bigr),
\end{aligned}
\]
in which the \(2k\)-ary \(T\) occurs only positively. Its operator
\(F_\psi\) on \(2k\)-tuples of the universe \(U\) of \(\mathbf A\oplus\mathbf
B\) has values in \(A^k\times B^k\), and \(T\) is consulted only at tuples in
\(A^k\times B^k\), so \(F_\psi(T)=\Gamma(T\cap(A^k\times B^k))\). Hence
\(\mathcal W=\Gamma(\mathcal W)=F_\psi(\mathcal W)\) is a fixed point of
\(F_\psi\), and every fixed point \(T'\) of \(F_\psi\) lies in
\(A^k\times B^k\) and satisfies \(\Gamma(T')=T'\), so \(\mathcal W\subseteq
T'\). Thus \(\operatorname{lfp}(F_\psi)=\mathcal W\), and by item~1 the
sentence
\[
\begin{aligned}
 (\exists\bar x)\Bigl(&\bigwedge_iD_1(x_i)\wedge{}\\
 &(\forall\bar y)\bigl(\bigvee_i\neg D_2(y_i)\vee[\mathbf{lfp}\,T\bar x\bar
 y\,.\,\psi](\bar x,\bar y)\bigr)\Bigr)
\end{aligned}
\]
holds in \(\mathbf A\oplus\mathbf B\) exactly when Duplicator does not win.

For finite structures let \(n\defeq|A|+|B|\). For fixed \(\sigma\) and \(k\)
the formula \(\vartheta\) is fixed, and \(F_\psi(T)\) is computed from \(T\)
by testing, for each of the at most \(n^{2k}\) tuples, \(\vartheta\) and the
at most \(kn^2\) membership questions of the second disjunct; this takes time
polynomial in \(n\). By \cref{thm:p03-w03-p03-059-finite-stages},
\(\operatorname{lfp}(F_\psi)\) is then computed in time polynomial in \(n\),
and the final test over \(A^k\times B^k\) adds at most \(n^{2k}\) lookups.

\emph{Item 3.} Let \(\mathbf A\) be a \(\sigma\)-structure. A family
\(\bar P=(P_{\bar b})_{\bar b\in B^k}\) of subsets of \(A^k\) corresponds to
\(\tau(\bar P)\defeq\{(\bar a,\bar b):\bar a\in P_{\bar b}\}\subseteq
A^k\times B^k\), and this correspondence is a bijection that preserves
inclusion in both directions. Let \(\Phi\) be the operator of
\(\pi_{\mathbf B}\) on \(\mathbf A\). We show that the \(T\)-components of
\(\Phi(\bar P)\) correspond to \(\Gamma(\tau(\bar P))\). A rule of the first
kind for \(\bar b\) puts \(\bar a\) into \(\Phi(\bar P)_{\bar b}\) exactly
when \(t^{\mathbf A}[\bar a]=t'^{\mathbf A}[\bar a]\) for terms with
\(t^{\mathbf B}[\bar b]\neq t'^{\mathbf B}[\bar b]\), that is, when
\(f_{\bar a\bar b}\) is not a function; a rule of the second kind does so
exactly when \(\mathbf A\models R(t_1,\ldots,t_r)\) at \(\bar a\) and
\(\mathbf B\not\models R(t_1,\ldots,t_r)\) at \(\bar b\). Together, as in item~2,
they put \(\bar a\) in exactly when \(f_{\bar a\bar b}\) is not a partial
homomorphism (with the rules for inequalities: not a one-to-one partial
homomorphism). A rule of the third kind for \(j\) fires under an assignment
\(\beta\) when \(\beta(y)\) satisfies
\((\beta(x_1),\ldots,\beta(y),\ldots,\beta(x_k))\in P_{\bar b[j/c]}\) for
every \(c\in B\), and \(\beta(x_j)\) is free; so it puts \(\bar a\) into
\(\Phi(\bar P)_{\bar b}\) exactly when some \(a\in A\) has
\((\bar a[j/a],\bar b[j/c])\in\tau(\bar P)\) for all \(c\in B\). This is
\(\Gamma\). The \(T\)-components of \(\Phi(\bar P)\) do not depend on the
goal component.

By induction on \(m\), the \(T\)-components of the stage \(\Phi^m\)
correspond to \(\Gamma^m(\emptyset)\), the \(m\)-th iterate of \(\Gamma\) from
the empty set. By the first paragraph of the proof of
\cref{thm:p03-w03-p03-061-datalog}, \(\Phi^\infty=\bigcup_m\Phi^m\) and
\(\Phi(\Phi^\infty)=\Phi^\infty\). So \(U\defeq\bigcup_m\Gamma^m(\emptyset)\)
corresponds to the \(T\)-components of \(\Phi^\infty\) and satisfies
\(\Gamma(U)=U\). Every fixed point \(T'\) of \(\Gamma\) contains each
\(\Gamma^m(\emptyset)\), by induction and monotonicity, hence contains \(U\);
so \(U=\mathcal W\). Finally, the goal rule fires in \(\Phi(\Phi^\infty)=
\Phi^\infty\) exactly when some \(\bar a\) lies in \(\Phi^\infty_{\bar b}\)
for all \(\bar b\in B^k\), that is, when some \(\bar a\) has \((\bar a,\bar
b)\in\mathcal W\) for all \(\bar b\). By item~1 this says that Duplicator does
not win.

Every rule body has at most \(k\) distinct variables: the terms of the first
two kinds use only \(x_1,\ldots,x_k\); the third kind uses \(y\) and the
\(x_i\) with \(i\neq j\); the goal body uses \(\bar x\). Heads have at most
\(k\) variables. So \(\pi_{\mathbf B}\) is a \(k\)-Datalog program, and a
\(k\)-Datalog\((\neq)\) program with the rules for inequalities.
\end{proof}
\end{annexitem}

\begin{annexitem}{annex:p03-w03-p03-061-horn-complete}{thm:p03-w03-p03-061-horn-complete}{A \(\mathsf P\)-complete existential game}
\begin{proof}
\emph{Membership.} The vocabulary, \(k=3\) and \(\mathbf B\) are fixed, so by
\cref{thm:p03-w03-p03-061-spoiler-lfp}, item~2, whether Duplicator wins the
\((\exists,3)\)-pebble game on \(\mathbf A\) and \(\mathbf B\) is decided in
time polynomial in \(|A|+2\).

\emph{Correctness of the reduction.} A Horn formula with at most three
literals per clause, after the normalization in
\cref{def:p03-w03-p03-061-horn}, is a finite list of clauses of its six
nonempty forms. By \cref{thm:p03-w03-p03-061-horn},
\(\varphi\) is satisfiable exactly when Duplicator wins the
\((\exists,3)\)-pebble game on \(\mathbf A_\varphi\) and \(\mathbf B\). A
formula without clauses is satisfiable and is sent to a fixed one-element
structure with empty relations, which maps homomorphically to \(\mathbf B\),
so Duplicator wins there by \cref{thm:p03-w03-p03-061-homomorphism}.

\emph{The reduction is computable in logarithmic space.} Let \(\varphi\) be
written as a list of clauses, each a list of at most three literals, a
literal being a sign followed by a variable name in binary. The machine keeps
a bounded number of pointers into the input, each of \(O(\log n)\) bits for
input length \(n\). A first scan detects an empty clause; if one occurs,
the machine outputs the fixed two-element structure whose \(P_1\) and
\(N_1\) both contain one element and whose other element is unused. This
is the encoding of the normalized contradiction. Otherwise it reserves a
separate output label for the fresh unused element and tags all input
variable names differently, avoiding a name collision without storing a
fresh long integer. To output the remaining universe it runs through the literal
positions \(p\) in order and writes the name at \(p\) when no earlier
position carries the same name, comparing names symbol by symbol with two
further pointers. To output the relations it runs through the clauses; for
each it counts the literals (at most three) and the positive ones (at most
one, as the clause is Horn), which fixes the symbol \(N_r\) or \(P_r\), and
writes the names of the negative literals in their order followed by the
name of the positive literal, if any; this is the tuple prescribed in
\cref{def:p03-w03-p03-061-horn}. Only the pointers and counters of bounded
number are stored.

\emph{Hardness.} Horn 3-satisfiability is \(\mathsf P\)-complete under
logarithmic-space reductions, as Kolaitis records with a reference
\cite[p.~76]{Kolaitis2003ExpressivePowerNotes}; this fact is used, not proved,
here. Let \(X\in\mathsf P\) and let \(r_1\) be a logarithmic-space reduction
of \(X\) to Horn 3-satisfiability and \(r_2\) the reduction above. The
composite \(r_2\circ r_1\) is computed in logarithmic space by running the
machine for \(r_2\) and, whenever it reads its \(i\)-th input symbol,
running the machine for \(r_1\) on \(x\) from the start while counting its
output symbols, discarding all but the \(i\)-th. The length of \(r_1(x)\) is
polynomial in \(|x|\), because a logarithmic-space machine has polynomially
many configurations, so the counter and the work space of \(r_2\) use
\(O(\log|x|)\) bits. Hence every problem in \(\mathsf P\) reduces in
logarithmic space to the game problem, which is in \(\mathsf P\): it is
\(\mathsf P\)-complete.
\end{proof}
\end{annexitem}

\begin{annexitem}{annex:p03-w03-p03-061-refines}{thm:p03-w03-p03-061-refines}{A win in the \(k\)-pebble game is a win in the existential games}
\begin{proof}
As \(\sigma\) has no constant symbols, \(C_{\mathbf A,\mathbf B}=\emptyset\).
Let \(\mathcal I\) be a winning family for the \(k\)-pebble game. A partial
isomorphism is an injective function on a finite set that preserves every
relation, hence a one-to-one partial homomorphism. For \(f\in\mathcal I\),
\(|f-\emptyset|=|f|\le k\); every subset of a member is a member; and the forth
property up to \(k\) is the same clause in both definitions. So \(\mathcal I\)
is a winning family for the \((\exists,\neq,k)\)-pebble game, and, as
one-to-one partial homomorphisms are partial homomorphisms, for the
\((\exists,k)\)-pebble game.

For the failure of the converse, the existential half is
\cref{thm:p03-w03-p03-061-cliques}. Suppose a winning family \(\mathcal I\)
for the \((k+1)\)-pebble game on \(\mathbf K_k\) and \(\mathbf K_{k+1}\)
existed, and let \(d_1,\ldots,d_{k+1}\) be the vertices of
\(\mathbf K_{k+1}\). Starting from \(\emptyset\in\mathcal I\), suppose
\(f_i\in\mathcal I\) has range \(\{d_1,\ldots,d_i\}\) and \(|f_i|=i\le k\).
The back property gives \(g\in\mathcal I\) with \(f_i\subseteq g\) and
\(d_{i+1}\) in its range; then \(f_{i+1}\defeq f_i\cup\{(g^{-1}(d_{i+1}),
d_{i+1})\}\subseteq g\) is in \(\mathcal I\) and has \(i+1\) pairs, since
\(d_{i+1}\) is not in the range of \(f_i\). The relation \(f_{k+1}\) is
injective with \(k+1\) pairs and domain in the \(k\)-element vertex set of
\(\mathbf K_k\), which is impossible. So Spoiler wins the
\((k+1)\)-pebble game.
\end{proof}
\end{annexitem}

\begin{annexitem}{annex:p03-w03-p03-061-switch}{thm:p03-w03-p03-061-switch}{Paths through a switch}
\begin{proof}
\emph{The ring.} The edges of \(\mathbf S\) between the twelve vertices other
than \(a,\ldots,h\) form the directed cycle
\[
 \alpha_1\to\alpha_2\to\alpha_3\to\alpha_4\to\beta'_1\to\beta'_2\to
 \alpha'_1\to\alpha'_2\to\alpha'_3\to\alpha'_4\to\beta_1\to\beta_2\to\alpha_1,
\]
the \emph{ring}; number its vertices \(0,\ldots,11\) in this order. Every
other edge joins an entry vertex to the ring or the ring to an exit vertex:
\(c\to0\), \(c\to6\), \(b\to2\), \(b\to8\), \(e\to4\), \(g\to10\), and
\(1\to h\), \(3\to a\), \(5\to d\), \(7\to f\), \(9\to a\), \(11\to d\). Each
of the twenty-four edges of \(\mathbf S\) is one of these. A ring vertex
therefore has as successors its ring successor and at most one exit vertex,
and the even-numbered ring vertices have no exit successor. A path passing
through \(\mathbf S\) consists of an entry vertex, a nonempty arc of
consecutive ring vertices \(i,i+1,\ldots\) (indices modulo \(12\)), and an
exit vertex.

\emph{Case \(Q\) starts at \(b\) and continues at \(2\).} Vertex \(2\) has
only the successor \(3\), so \(3\in Q\). The successors of \(3\) are \(a\)
and \(4\); \(a\) lies on \(P\), so \(4\in Q\), and \(4\) has only the successor
\(5\), so \(5\in Q\). The vertex before \(a\) on \(P\) is \(3\) or \(9\); as
\(3\in Q\), it is \(9\). The entry vertex of \(P\) is not \(b\in Q\). If \(P\)
entered at \(g\) (ring vertex \(10\)) or at \(c\) through \(0\), its arc,
running forward to \(9\), would contain \(2\in Q\); if it entered at \(e\), it
would start at \(4\in Q\). So \(P\) enters at \(c\) through \(6\) and is
\(c,6,7,8,9,a=q(c,a)\). The successors of \(5\) are \(d\) and \(6\in P\), so
\(Q=b,2,3,4,5,d=q(b,d)\).

Let \(R\) be a path passing through \(\mathbf S\) with no vertex in common with
\(P\) or \(Q\). Its entry is not \(c\) or \(b\), nor \(e\), whose only
successor \(4\) lies on \(Q\); so it is \(g\), followed by \(10\) and \(11\).
The successors of \(11\) are \(d\in Q\) and \(0\); then \(0\) has only the
successor \(1\); the successors of \(1\) are \(h\) and \(2\in Q\). So
\(R=g,10,11,0,1,h=q(g,h)\), which indeed avoids \(P\) and \(Q\).

\emph{Case \(Q\) continues at \(8\).} The mirror \(\mu\) is an automorphism of
\(\mathbf S\) fixing \(a,b,c,d\) and moving each ring vertex \(i\) to
\(i+6\); it carries \(8\) to \(2\). Applying the first case to \(\mu(P)\)
and \(\mu(Q)\) gives \(\mu(P)=q(c,a)\), \(\mu(Q)=q(b,d)\) and the unique third
path \(q(g,h)\); applying \(\mu\) again gives \(P=p(c,a)\), \(Q=p(b,d)\) and the
unique third path \(p(e,f)\).

These are the only cases, since the successors of \(b\) are \(2\) and \(8\).
\end{proof}
\end{annexitem}

\begin{annexitem}{annex:p03-w03-p03-061-templates}{thm:p03-w03-p03-061-templates}{Templates are walks that are injective on good positions}
\begin{proof}
\emph{Item 1.} Inside each block \(\mathrm{CA}_\nu(o)\), \(\mathrm{BD}_\nu(o)\),
\(q(g,h)\langle o\rangle\) and \(p(e,f)\langle o\rangle\), consecutive entries
are joined by edges of the switch of \(o\), copied by item~1 of
\cref{def:p03-w03-p03-061-fhw-graph}. The junctions are edges by the other
items: \(z_1\to c\langle o_N\rangle\), \(a\langle o_{r+1}\rangle\to
c\langle o_r\rangle\) and \(a\langle o_1\rangle\to z_2\) in \(T_1\);
\(z_3\to b\langle o_1\rangle\), \(d\langle o_r\rangle\to b\langle
o_{r+1}\rangle\), \(d\langle o_N\rangle\to u_1\); within \(\Gamma_i\), the
edges \(u_i\to g\langle o^y_1\rangle\), \(h\langle o^y_s\rangle\to
g\langle o^y_{s+1}\rangle\), \(h\langle o^y_{m(y)}\rangle\to v_i\), or
\(u_i\to v_i\) when \(m(y)=0\); \(v_i\to u_{i+1}\), \(v_k\to n_0\);
\(n_{j-1}\to e\langle\kappa(j)\rangle\) and \(f\langle\kappa(j)\rangle\to
n_j\), as \(\kappa(j)\) is an occurrence in clause \(j\); and \(n_l\to z_4\).

\emph{Item 2.} Vertices in different switches are different, and each of
\(z_1,\ldots,z_4\), \(u_i\), \(v_i\), \(n_j\) occurs at exactly one position
of \(T_1(\nu)T_2(\nu,\kappa)\). Within one block the entries are distinct.
Fix an occurrence \(o\) with literal \(y\); the entries in its switch at the
considered positions come from: \(\mathrm{CA}_\nu(o)\) and
\(\mathrm{BD}_\nu(o)\); \(q(g,h)\langle o\rangle\) when \(y\) is the literal
\(\lambda^\nu_i\) of its variable, that is, when \(\nu\) makes \(y\) false;
and \(p(e,f)\langle o\rangle\) when \(o=\kappa(j)\) for the clause \(j\) of
\(o\) and the segment of \(j\) meets \(D\), in which case \(\nu\) makes \(y\)
true because \(D\) is good. If \(\nu\) makes \(y\) true, the mode is \(p\) and
the blocks present are among \(p(c,a),p(b,d),p(e,f)\); if it makes \(y\)
false, the mode is \(q\) and they are among \(q(c,a),q(b,d),q(g,h)\). In both
cases they are pairwise vertex-disjoint by
\cref{def:p03-w03-p03-061-switch}, and each occurs once, since \(o\) lies in
one column and in one clause.

\emph{Item 3.} \(\mathrm{CA}_\nu(o_r)\) and \(\mathrm{BD}_\nu(o_r)\) have six
entries in either mode, so the positions of these blocks in \(T_1\) and at
the start of \(T_2\) do not depend on \(\nu\). Their first and last entries are
\(c,a\) and \(b,d\) of the switch of \(o_r\) in either mode, and their inner
entries depend on the mode, which depends only on \(\nu\) at the variable of
\(o_r\). In a balanced formula \(\Gamma_i\) has \(6m(x_i)+2\) entries whichever
literal \(\lambda^\nu_i\) is, so the positions of \(u_i\), \(v_i\) and of all
later entries do not depend on \(\nu\); the entries strictly between \(u_i\)
and \(v_i\) are determined by \(\lambda^\nu_i\), that is, by \(\nu(x_i)\). The
segment of clause \(j\) holds \(p(e,f)\langle\kappa(j)\rangle\), determined
by \(\kappa(j)\). The remaining entries are \(z_1,\ldots,z_4\), the vertices
\(c,a,b,d\) of the switches, \(u_i,v_i\) and \(n_j\), at fixed positions.
\end{proof}
\end{annexitem}

\begin{annexitem}{annex:p03-w03-p03-061-fhw}{thm:p03-w03-p03-061-fhw}{The reduction of Fortune, Hopcroft and Wyllie}
\begin{proof}
\emph{Entering and leaving switches.} By the list of edges in
\cref{def:p03-w03-p03-061-fhw-graph}, every edge of \(\mathbf G_\varphi\) that
enters the switch of an occurrence \(o\) from outside ends at
\(b\langle o\rangle\), \(c\langle o\rangle\), \(e\langle o\rangle\) or
\(g\langle o\rangle\), and every edge that leaves it starts at
\(a\langle o\rangle\), \(d\langle o\rangle\), \(f\langle o\rangle\) or
\(h\langle o\rangle\). Let \(P\) be a simple path whose first and last
vertices lie outside the switch of \(o\). A maximal run of consecutive
vertices of \(P\) inside that switch is entered by an edge from outside and
left by an edge to outside, so it is a path of the switch from an entry to an
exit vertex: it passes through the switch. The vertices \(z_1,z_3\) have no
incoming and \(z_2,z_4\) no outgoing edge.

\emph{Paths give a satisfying assignment.} Let \(P_1\) from \(z_1\) to
\(z_2\) and \(P_2\) from \(z_3\) to \(z_4\) be simple paths with no common
vertex. We show by induction on \(r=1,\ldots,N\): the last
\(6r+1\) vertices of \(P_1\) are \(\mathrm{CA}(o_r),\ldots,\mathrm{CA}(o_1),z_2\)
and the first \(6r+1\) vertices of \(P_2\) are
\(z_3,\mathrm{BD}(o_1),\ldots,\mathrm{BD}(o_r)\), where for each \(o_s\) the
two blocks are either \(p(c,a)\langle o_s\rangle,p(b,d)\langle o_s\rangle\)
or \(q(c,a)\langle o_s\rangle,q(b,d)\langle o_s\rangle\); call \(o_s\) of
\emph{mode} \(p\) or \(q\) accordingly. The vertex of \(P_1\) before
\(z_2\), respectively before \(c\langle o_{r-1}\rangle\) when \(r\ge2\), is the
only in-neighbour \(a\langle o_r\rangle\) of that vertex; the vertex of \(P_2\)
after \(z_3\), respectively after \(d\langle o_{r-1}\rangle\), is the only
out-neighbour \(b\langle o_r\rangle\). The run of \(P_1\) in the switch of
\(o_r\) that ends at \(a\langle o_r\rangle\) and the run of \(P_2\) that starts
at \(b\langle o_r\rangle\) pass through the switch and have no common vertex,
so \cref{thm:p03-w03-p03-061-switch} identifies them as stated. This
completes the induction. The only in-neighbour of \(c\langle o_N\rangle\) is
\(z_1\), the first vertex of \(P_1\), so \(P_1\) is \(z_1,\mathrm{CA}(o_N),
\ldots,\mathrm{CA}(o_1),z_2\); the only out-neighbour of \(d\langle o_N\rangle\)
is \(u_1\), and \(P_2\) continues from \(u_1\) by a path \(R\) to \(z_4\).

\(R\) has no vertex in common with \(P_1\) or with the first part of
\(P_2\). If \(R\) meets the switch of \(o\), each of its runs there passes
through the switch avoiding the two blocks of \(P_1\) and \(P_2\), so by
\cref{thm:p03-w03-p03-061-switch} it is \(p(e,f)\langle o\rangle\) when \(o\)
has mode \(p\), and \(q(g,h)\langle o\rangle\) when \(o\) has mode \(q\).
Follow \(R\) from \(u_1\). From \(u_i\) it moves to \(v_i\) or to
\(g\langle o^y_1\rangle\) for a literal \(y\) of \(x_i\); entering at
\(g\) it runs through \(q(g,h)\), so the switch has mode \(q\), and leaves at
\(h\), whose only out-neighbour is \(g\) of the next occurrence of \(y\) or
\(v_i\). So \(R\) runs through the whole column of one literal \(y_i\) of
\(x_i\), all of whose occurrences have mode \(q\), and reaches \(v_i\), whose
only out-neighbour is \(u_{i+1}\), respectively \(n_0\). From \(n_{j-1}\) it
moves to \(e\langle o\rangle\) for an occurrence \(o\) in clause \(j\), runs
through \(p(e,f)\langle o\rangle\), so \(o\) has mode \(p\), and reaches
\(n_j\). Let \(\nu\) make each \(y_i\) false. For clause \(j\) with the
occurrence \(o=(y,j)\) used by \(R\): if \(\nu\) made \(y\) false, then
\(y=y_i\) for its variable \(x_i\), and every occurrence of \(y\) would have
mode \(q\), against the mode \(p\) of \(o\). So \(\nu\) makes \(y\) true, and
\(\nu\) satisfies every clause.

\emph{A satisfying assignment gives paths.} Let \(\nu\) satisfy \(\varphi\)
and let \(\kappa(j)\) be an occurrence in clause \(j\) of a literal made true
by \(\nu\). The set \(D\) of all positions of \(T_2\) is good for
\((\nu,\kappa)\). By \cref{thm:p03-w03-p03-061-templates}, items~1 and~2,
\(T_1(\nu)\) and \(T_2(\nu,\kappa)\) are walks without repeated vertices, so
simple paths, from \(z_1\) to \(z_2\) and from \(z_3\) to \(z_4\), with no
common vertex.

\emph{The complete formulas.} \(\varphi_k\) is unsatisfiable: under any
assignment, the clause consisting of the literals it makes false is false.
\end{proof}
\end{annexitem}

\begin{annexitem}{annex:p03-w03-p03-061-family}{thm:p03-w03-p03-061-family}{Duplicator's winning family}
\begin{proof}
\emph{The two queries.} The two paths of \(\mathbf A_k\) have no common
vertex and join \(d_{s_1}\) to \(d_{s_2}\) and \(d_{s_3}\) to \(d_{s_4}\); as
\(L_1,L_2\ge2\), the four constants are distinct, so \(\mathbf A_k\in
\mathcal F_{\mathbf H_1}\) and \(\mathrm{FiSH}(\mathbf H_1)(\mathbf A_k)=1\).
\(\mathbf B_k\in\mathcal F_{\mathbf H_1}\), and
\(\mathrm{FiSH}(\mathbf H_1)(\mathbf B_k)=0\) by
\cref{thm:p03-w03-p03-061-fhw}.

\emph{The formula \(\varphi_k\) is balanced.} Each literal lies in exactly
the \(2^{k-1}\) clauses that choose it for its variable, so
\(m(x_i)=m(\bar x_i)\). Hence \(L_1,L_2\) and the positions in
\cref{thm:p03-w03-p03-061-templates}, item~3, are fixed, and
\(h_{\nu,\kappa}\) is defined on all of \(A_k\).

\emph{Members are one-to-one partial homomorphisms.} Let
\(f=h_{\nu,\kappa}\restriction(D\cup W)\) with \(D\) good. The constants go
to the constants: \(\eta_0\mapsto z_1\), \(\eta_{L_1-1}\mapsto z_2\),
\(\eta'_0\mapsto z_3\), \(\eta'_{L_2-1}\mapsto z_4\), the first and last
entries of the templates; so \(C\subseteq f\), where \(C\) is the set of these
four pairs. An edge of \(\mathbf A_k\) is \((\eta_t,\eta_{t+1})\) or
\((\eta'_t,\eta'_{t+1})\), and goes to an edge by item~1 of
\cref{thm:p03-w03-p03-061-templates}. The positions of \(\eta'_0\) and
\(\eta'_{L_2-1}\) lie outside every clause segment, so item~2 shows that
\(f\) is injective. Moreover \(|f-C|=|D|\le k\).

\emph{Nonempty and closed under subfunctions.} Clauses are nonempty, so
clause choices exist, and \(D=\emptyset\) is good: \(h_{\nu,\kappa}
\restriction W=C\in\mathcal I_k\). If \(C\subseteq g\subseteq f\), then
\(g=h_{\nu,\kappa}\restriction(D'\cup W)\) for some \(D'\subseteq D\), and a
subset of a good set is good.

\emph{Forth.} Let \(f=h_{\nu,\kappa}\restriction(D\cup W)\in\mathcal I_k\)
with \(|D|<k\), and \(a\in A_k\). If \(a\in D\cup W\) there is nothing to
do. Let \(F\) be the set of variables \(x_i\) such that some vertex of \(D\)
sits at a position inside a block \(\mathrm{CA}(o_r)\) or \(\mathrm{BD}(o_r)\)
other than its first and last, with \(x_i\) the variable of \(o_r\); or
strictly between the positions of \(u_i\) and \(v_i\); or in the segment of a
clause \(j\) with \(x_i\) the variable of the literal of \(\kappa(j)\). Let
\(K\) be the set of clauses whose segment meets \(D\). Each vertex of \(D\)
contributes at most one variable, so \(|F|\le|D|\le k-1\). By item~3 of
\cref{thm:p03-w03-p03-061-templates}, if \(\nu'\) agrees with \(\nu\) on
\(F\) and \(\kappa'\) with \(\kappa\) on \(K\), then \(h_{\nu',\kappa'}\) and
\(h_{\nu,\kappa}\) agree on \(D\cup W\), and \(D\) is good for
\((\nu',\kappa')\), because for \(j\in K\) the literal of
\(\kappa'(j)=\kappa(j)\) has its variable in \(F\).

If \(a\) does not lie in the segment of a clause outside \(K\), the set
\(D\cup\{a\}\) is good for \((\nu,\kappa)\), and
\(h_{\nu,\kappa}\restriction(D\cup\{a\}\cup W)\in\mathcal I_k\) extends
\(f\). Otherwise \(a\) lies in the segment of a clause \(j\notin K\). That
clause contains a literal of each of the \(k\) variables, and \(|F|<k\), so it
contains a literal \(y\) of a variable \(x_i\notin F\). Let \(\nu'\) agree
with \(\nu\) except that it makes \(y\) true, and \(\kappa'\) with \(\kappa\)
except that \(\kappa'(j)\defeq(y,j)\). Then \(\nu'\) agrees with \(\nu\) on
\(F\) and \(\kappa'\) with \(\kappa\) on \(K\), so
\(g\defeq h_{\nu',\kappa'}\restriction(D\cup\{a\}\cup W)\) extends \(f\);
\(D\cup\{a\}\) is good for \((\nu',\kappa')\), since the clauses of \(K\)
keep their true literals and clause \(j\) receives one; and
\(|D\cup\{a\}|\le k\). So \(g\in\mathcal I_k\) and \(a\in\operatorname{dom}(g)\).
\end{proof}
\end{annexitem}

\begin{annexitem}{annex:p03-w03-p03-061-h1}{thm:p03-w03-p03-061-h1}{The 2-disjoint paths query is not existentially definable}
\begin{proof}
Suppose a sentence \(\chi\) of \(\exists\mathrm L^k_{\infty\omega}(\neq)\)
defined \(\mathrm{FiSH}(\mathbf H_1)\) on \(\mathcal F_{\mathbf H_1}\). By
\cref{thm:p03-w03-p03-061-family} and \cref{thm:p03-w03-p03-061-strategy},
Duplicator wins the \((\exists,\neq,k)\)-pebble game on \(\mathbf A_k\) and
\(\mathbf B_k\), so by \cref{thm:p03-w03-p03-061-game-logic} every
\(\exists\mathrm L^k_{\infty\omega}(\neq)\)-sentence true in \(\mathbf A_k\)
is true in \(\mathbf B_k\). As \(\mathbf A_k\models\chi\), also
\(\mathbf B_k\models\chi\), contradicting
\(\mathrm{FiSH}(\mathbf H_1)(\mathbf B_k)=0\). A sentence of
\(\exists\mathrm L^\omega_{\infty\omega}(\neq)\) lies in some
\(\exists\mathrm L^k_{\infty\omega}(\neq)\). A Datalog\((\neq)\) program is a
\(k\)-Datalog\((\neq)\) program for \(k\) the largest number of variables of
a rule, and its Boolean goal is then
\(\exists\mathrm L^k_{\infty\omega}(\neq)\)-definable by
\cref{thm:p03-w03-p03-061-datalog}.
\end{proof}
\end{annexitem}

\begin{annexitem}{annex:p03-w03-p03-061-quotient}{thm:p03-w03-p03-061-quotient}{Identifying distinguished vertices}
\begin{proof}
Write \(C\) for \(\{(a_i,b_i):i\le n\}\) and \(C'\) for
\(\{([a_i],[b_i]):i\le n\}\), the relations of the constants. A class \([x]\)
with more than one element consists of distinguished vertices. Let
\(f\in\mathcal I\); then \(C\subseteq f\).

\(f/{\equiv}\) is a function: if \([x]=[x']\) with \(x\neq x'\), then
\(x=a_i\), \(x'=a_j\) with \(i\equiv j\), so \(f(x)=b_i\), \(f(x')=b_j\) and
\([b_i]=[b_j]\). It is injective: if \([f(x)]=[f(x')]\) with \(f(x)\neq
f(x')\), then \(f(x)=b_i\), \(f(x')=b_j\) with \(i\equiv j\); as
\(f(a_i)=b_i\) and \(f\) is injective, \(x=a_i\), \(x'=a_j\), so
\([x]=[x']\). It preserves edges: an edge \(([x],[x'])\) of \(\mathbf
A/{\equiv}\) between classes in its domain comes from an edge \(x_0\to x_0'\)
of \(\mathbf A\) with \(x_0\in[x]\), \(x'_0\in[x']\); every member of a class
in the domain of \(f/{\equiv}\) lies in the domain of \(f\) (a non-singleton
class consists of constants), so \(f(x_0)\to f(x'_0)\) in \(\mathbf B\) and
\(([f(x_0)],[f(x'_0)])\) is an edge of \(\mathbf B/{\equiv}\). And
\(C'\subseteq f/{\equiv}\).

The map \((x,y)\mapsto([x],[y])\) sends \(f-C\) onto \(f/{\equiv}-C'\): a
pair of \(C\) goes into \(C'\), and a pair \((x,y)\) of \(f-C\) has \(x\) not
distinguished (as \(f(a_i)=b_i\)), so \([x]=\{x\}\) is not the class of a
constant and \(([x],[y])\notin C'\). It is injective on \(f-C\), since
\(x\mapsto[x]\) is injective on non-distinguished vertices. Hence
\(|f/{\equiv}-C'|=|f-C|\le k\).

Closure under subfunctions: let \(C'\subseteq g'\subseteq f/{\equiv}\) and
\(g\defeq C\cup\{(x,y)\in f:([x],[y])\in g'\}\). Then \(C\subseteq g\subseteq
f\), so \(g\in\mathcal I\), and \(g/{\equiv}=g'\): the inclusion
\(\subseteq\) holds as \(C'\subseteq g'\), and each pair of \(g'\) is
\(([x],[y])\) for some \((x,y)\in f\), which then lies in \(g\).

Forth: let \(|f/{\equiv}-C'|<k\) and let \([x]\) be a vertex of
\(\mathbf A/{\equiv}\). Then \(|f-C|<k\), and the forth property of
\(\mathcal I\) gives \(g\in\mathcal I\) with \(f\subseteq g\) and \(x\in
\operatorname{dom}(g)\); so \(f/{\equiv}\subseteq g/{\equiv}\) and \([x]\) is
in the domain of \(g/{\equiv}\). The family is nonempty since \(\mathcal I\)
is.
\end{proof}
\end{annexitem}

\begin{annexitem}{annex:p03-w03-p03-061-h2h3}{thm:p03-w03-p03-061-h2h3}{The basic patterns \(\mathbf H_2\) and \(\mathbf H_3\)}
\begin{proof}
Fix \(k\). In \(\mathbf A_k\) and \(\mathbf B_k\) the distinguished vertices
are, in order,
\[
 (\eta_0,\eta_{L_1-1},\eta'_0,\eta'_{L_2-1})
 \quad\hbox{and}\quad(z_1,z_2,z_3,z_4).
\] In both graphs the first and third have no incoming and
the second and fourth no outgoing edges. If a structure has several constant
symbols denoting the same element and all but one of them are dropped, a
winning family for the original structures is one for the reducts, because
the set of pairs of constants is unchanged.

\emph{\(\mathbf H_2\).} Identify the second and third distinguished
vertices: \(\mathbf A^2_k\defeq\mathbf A_k/{\equiv}\) and \(\mathbf
B^2_k\defeq\mathbf B_k/{\equiv}\) with \(2\equiv3\), read as
\(\sigma_{\mathbf H_2}\)-structures with \(d_{s_1},d_{s_2},d_{s_3}\) denoting
the classes of the first, second (\(=\)third) and fourth vertex. These are
distinct, so both structures lie in \(\mathcal F_{\mathbf H_2}\), and by
\cref{thm:p03-w03-p03-061-family,thm:p03-w03-p03-061-quotient} Duplicator
wins the \((\exists,\neq,k)\)-pebble game on them. In \(\mathbf A^2_k\) the
paths \(\eta_0,\ldots,\eta_{L_1-1}\) and \(\eta'_0,\ldots,\eta'_{L_2-1}\)
become simple paths from \(d_{s_1}\) to \(d_{s_2}\) and from \(d_{s_2}\) to
\(d_{s_3}\) sharing only \(d_{s_2}\), so \(\mathrm{FiSH}(\mathbf
H_2)(\mathbf A^2_k)=1\). Suppose simple paths \(P\) from \([z_1]\) to
\([z_2]=[z_3]\) and \(Q\) from \([z_2]\) to \([z_4]\) in \(\mathbf B^2_k\)
share only \([z_2]\). The last edge of \(P\) enters \([z_2]\), so it comes
from an edge of \(\mathbf B_k\) into \(z_2\) (\(z_3\) has none); the first edge
of \(Q\) comes from an edge out of \(z_3\) (\(z_2\) has none). All other
vertices of \(P\) and \(Q\) are singletons, and edges between them are edges
of \(\mathbf B_k\). Replacing \([z_2]\) by \(z_2\) in \(P\) and by \(z_3\) in
\(Q\) gives simple paths of \(\mathbf B_k\) from \(z_1\) to \(z_2\) and from
\(z_3\) to \(z_4\) with no common vertex, contradicting
\(\mathrm{FiSH}(\mathbf H_1)(\mathbf B_k)=0\). So
\(\mathrm{FiSH}(\mathbf H_2)(\mathbf B^2_k)=0\).

\emph{\(\mathbf H_3\).} Identify the first with the fourth and the second with
the third distinguished vertex, and let \(d_{s_1},d_{s_2}\) denote the classes
\([z_1]=[z_4]\) and \([z_2]=[z_3]\) (in \(\mathbf A_k\) the corresponding
classes). They are distinct, the structures lie in \(\mathcal F_{\mathbf
H_3}\), and Duplicator wins the \((\exists,\neq,k)\)-pebble game on them by
the same two results. In the quotient of \(\mathbf A_k\) the two paths become
simple paths from \(d_{s_1}\) to \(d_{s_2}\) and from \(d_{s_2}\) to
\(d_{s_1}\) that share only their endpoints, so the query holds. Suppose
simple paths \(P\) from \([z_1]\) to \([z_2]\) and \(Q\) from \([z_2]\) to
\([z_1]\) in the quotient of \(\mathbf B_k\) share only their endpoints. The
first edge of \(P\) leaves \([z_1]\), hence comes from an edge out of \(z_1\)
(\(z_4\) has none); its last edge comes from an edge into \(z_2\). The first
edge of \(Q\) comes from an edge out of \(z_3\), its last from an edge into
\(z_4\) (\(z_1\) has none). Simple paths visit their endpoints only at the
ends, so replacing the endpoints accordingly gives simple paths of
\(\mathbf B_k\) from \(z_1\) to \(z_2\) and from \(z_3\) to \(z_4\) with no
common vertex, a contradiction as before.

In both cases, if an \(\exists\mathrm L^k_{\infty\omega}(\neq)\)-sentence
defined the query, it would hold in the positive structure and, by
\cref{thm:p03-w03-p03-061-game-logic}, in the negative one.
\end{proof}
\end{annexitem}

\begin{annexitem}{annex:p03-w03-p03-061-single-loop}{thm:p03-w03-p03-061-single-loop}{The single-loop boundary example}
\begin{proof}
At stage \(n\ge1\), the unary relation computed by the program consists
exactly of the vertices reachable from \(c\) by a directed walk of length
between one and \(n\). This follows by induction: the first rule supplies
the walks of length one and the second appends one edge. Hence the least
fixed point contains \(c\) precisely when there is a positive-length
closed walk from \(c\) to itself.

A simple directed cycle through \(c\) is such a walk. Conversely choose
a closed walk from \(c\) to itself of least positive length. No internal
vertex is \(c\), since its first return would be a shorter positive closed
walk. No two internal vertices coincide, since deleting the segment
between equal vertices would give a shorter positive closed walk with the
same endpoints. Thus the chosen walk is a simple directed cycle, including
the length-one case.

For the finite-variable claim, unwind the stages. Write \(\rho_1(x)=E(c,x)\)
and recursively
\[
 \rho_{n+1}(x)=E(c,x)\vee\exists y\,(\rho_n(y)\wedge E(y,x)).
\]
Only \(x,y\) are needed: in the copy of \(\rho_n(y)\), bound occurrences
of \(x\) are confined to that subformula and do not bind the \(x\) in
\(E(y,x)\). The resulting sentence
\(\bigvee_{n\ge1}\rho_n(c)\) belongs to
\(\exists\mathrm L^2_{\infty\omega}\) and defines the same query.
Finally, a loop has coincident endpoints, whereas each edge in the star
definition joins its centre to a distinct leaf. It is therefore a
non-star pattern with a definable query under the stated extension.
\end{proof}
\end{annexitem}

\begin{annexitem}{annex:p03-w03-p03-061-dichotomy}{thm:p03-w03-p03-061-dichotomy}{Descriptive dichotomy for fixed subgraph homeomorphism queries}
\begin{proof}
\emph{Item 1} is \cref{thm:p03-w03-p03-061-star} together with
\cref{thm:p03-w03-p03-061-star-datalog}.

\emph{Item 2.} By \cref{thm:p03-w03-p03-061-h1,thm:p03-w03-p03-061-h2h3},
for every \(k\) none of \(\mathrm{FiSH}(\mathbf H_1)\),
\(\mathrm{FiSH}(\mathbf H_2)\), \(\mathrm{FiSH}(\mathbf H_3)\) is
\(\exists\mathrm L^k_{\infty\omega}(\neq)\)-definable. Let \(\mathbf H\) be
a pattern graph that is not a star graph. By
\cref{thm:p03-w03-p03-061-basic} some \(\mathbf H_i\) is isomorphic to a
subgraph of \(\mathbf H\), and \(\mathbf H_i\) is itself a pattern graph. By
\cref{thm:p03-w03-p03-061-transfer}, for every \(k\),
\(\mathrm{FiSH}(\mathbf H)\) is not
\(\exists\mathrm L^k_{\infty\omega}(\neq)\)-definable on
\(\mathcal F_{\mathbf H}\). A sentence of
\(\exists\mathrm L^\omega_{\infty\omega}(\neq)\) lies in one of these logics,
and a Datalog\((\neq)\) program defines an
\(\exists\mathrm L^k_{\infty\omega}(\neq)\)-definable query for some \(k\),
as in the proof of \cref{thm:p03-w03-p03-061-h1}.
\end{proof}
\end{annexitem}
